\documentclass[11pt]{article}

\usepackage[review]{acl}

\usepackage{times}
\usepackage{latexsym}
\usepackage[T1]{fontenc}
\usepackage[utf8]{inputenc}
\usepackage{microtype}
\usepackage{inconsolata}

\usepackage{booktabs}
\usepackage{amsfonts}
\usepackage{amsmath}
\usepackage{amssymb}
\usepackage{nicefrac}
\usepackage{xcolor}
\usepackage{graphicx}
\usepackage{algorithm}
\usepackage{algorithmic}
\usepackage{amsthm}
\usepackage{mathtools}
\usepackage{multirow}
\usepackage{enumitem}
\usepackage{subcaption}

\newcommand{\gV}{\mathcal{V}}
\newcommand{\gS}{\mathcal{S}}
\newcommand{\gA}{\mathcal{A}}
\newcommand{\gD}{\mathcal{D}}
\newcommand{\gL}{\mathcal{L}}
\newcommand{\gO}{\mathcal{O}}
\newcommand{\gN}{\mathcal{N}}
\newcommand{\rvx}{\mathbf{x}}
\newcommand{\rvy}{\mathbf{y}}
\newcommand{\rvz}{\mathbf{z}}
\newcommand{\rvi}{\mathbf{i}}
\newcommand{\rvm}{\mathbf{m}}
\newcommand{\rvh}{\mathbf{h}}
\newcommand{\rvw}{\mathbf{w}}
\newcommand{\ind}{\mathbf{1}}
\newcommand{\E}{\mathbb{E}}
\newcommand{\R}{\mathbb{R}}
\newcommand{\softmax}{\text{softmax}}
\newcommand{\sigmoid}{\sigma}
\DeclareMathOperator*{\argmax}{arg\,max}
\DeclareMathOperator*{\argmin}{arg\,min}

\newtheorem{theorem}{Theorem}[section]
\newtheorem{proposition}[theorem]{Proposition}
\newtheorem{definition}[theorem]{Definition}
\newtheorem{assumption}[theorem]{Assumption}
\newtheorem{corollary}[theorem]{Corollary}
\newtheorem{lemma}[theorem]{Lemma}

\theoremstyle{remark}
\newtheorem{remark}[theorem]{Remark}

\title{Not All Tokens Are Equal: Importance-Aware Masking\\for Discrete Diffusion Language Models}

\author{Adele Chinda \\
  Georgia State University \\
  \texttt{achinda1@gsu.edu} \\\And
  Haoxin Wang \\
  Georgia State University \\
  \texttt{haoxinwang@gsu.edu} \\}

\begin{document}
\maketitle

\begin{abstract}
Discrete diffusion language models (DLMs) offer an appealing alternative to autoregressive (AR) generation by enabling parallel decoding and bidirectional context, yet they consistently underperform AR baselines in perplexity.
A fundamental limitation is their uniform treatment of all tokens during the masking process: content-bearing words and function words are masked with identical probability, wasting model capacity on trivial predictions.
We introduce \textbf{ATAT} (Adaptive Token Attention for Text Diffusion), a framework that learns token-level importance through a hybrid estimator combining contextual signals from a frozen GPT-2 anchor with corpus frequency priors ($i^l = 0.7 \cdot i^l_{\text{context}} + 0.3 \cdot i^l_{\text{freq}}$).
The hybrid achieves 0.82 correlation with oracle surprisal and 39.03 perplexity on WikiText-2, outperforming both contextual-only (40.12) and frequency-only (41.87) ablations.
\end{abstract}

%==============================================================================
\section{Introduction}
\label{sec:intro}
%==============================================================================

Language modeling has been dominated by the autoregressive (AR) paradigm, where each token is generated conditioned on all preceding tokens~\citep{gpt2,gpt3,llama,llama2}. While remarkably effective, this left-to-right constraint prevents parallel generation and limits the model's ability to exploit bidirectional context during inference. \textbf{Discrete diffusion language models} (DLMs)~\citep{d3pm,mdlm,sedd} offer a principled alternative by framing text generation as an iterative denoising process: a clean sequence $\rvx$ is progressively corrupted by masking tokens, and a neural network learns to reverse this corruption. The forward (corruption) process follows:
\begin{equation}
    q(\rvz_t \mid \rvx) = \prod_{l=1}^{L} \bigl[\alpha_t \cdot \ind[z_t^l = x^l] + (1{-}\alpha_t) \cdot \ind[z_t^l = \texttt{[M]}]\bigr],
    \label{eq:forward}
\end{equation}
where $\alpha_t$ monotonically decreases from 1 (clean data) to 0 (fully masked), and $\texttt{[M]}$ is an absorbing mask token. At generation time, the model begins from fully masked input and iteratively predicts tokens at masked positions using bidirectional attention, enabling parallel decoding across all positions simultaneously.

Despite these architectural advantages, DLMs consistently lag behind comparably-sized AR models in perplexity and downstream generation quality~\citep{mdlm,sedd,adlm}. Recent work has narrowed this gap through improved noise schedules~\citep{mdlm}, score-entropy objectives~\citep{sedd}, and frequency-based anchor tokens~\citep{adlm}, yet a fundamental limitation persists: \textit{all existing DLMs treat every token identically during masking}.

\paragraph{The Uniform Masking Problem.} In Eq.~\ref{eq:forward}, every position $l$ has the same masking probability $(1{-}\alpha_t)$, regardless of the token's informational content. This is linguistically naive. Natural language exhibits massive heterogeneity in token informativeness: content words like ``quantum'' or ``discover'' carry 8--12 bits of surprisal, while function words like ``the'' or ``of'' carry merely 2--4 bits~\citep{zipf1949human,shannon1948mathematical,shannon1951prediction}. Uniform masking forces the model to allocate equal capacity to reconstructing both trivial function words (solvable from local syntax) and rare content words (requiring world knowledge and long-range dependencies). This is analogous to training a student by testing equally on material they already know and material they have never seen---a highly inefficient allocation of learning signal.

\paragraph{Limitations of Frequency-Based Anchoring.} ADLM~\citep{adlm} partially addresses this by identifying ``anchor tokens'' based on corpus frequency, preserving rare tokens during the diffusion process. However, frequency is a coarse proxy for importance. A rare proper noun and a rare technical term have identical frequency but potentially very different contextual significance. The word ``bank'' has the same frequency whether it anchors a financial discussion (high importance) or appears peripherally in nature writing (low importance). Frequency-based methods also systematically undervalue common tokens that are contextually critical---for instance, ``the'' in ``\textit{The one who survived}'' is grammatically common but semantically essential.

\begin{figure*}[t]
    \centering
    \includegraphics[width=.7\textwidth]{figures/atat_architecture.pdf}
    \caption{Importance-aware denoising in ATAT. Starting from fully masked tokens ($t{=}1$), the model progressively reveals words guided by learned importance scores $\rvi$. High-importance tokens (e.g., semantically rich content words, shown in red) are unmasked first, providing contextual anchors that guide the recovery of medium-importance (yellow) and low-importance (green) tokens in later steps. By $t{=}0$ all tokens are recovered. This importance-driven schedule contrasts with uniform random masking, allowing the model to build meaning hierarchically from the most informative tokens outward.}
    \label{fig:architecture}
\end{figure*}

\paragraph{Our Approach: Learned Hybrid Importance.} We ask: \textit{can we learn which tokens matter and adaptively mask them?} We introduce \textbf{ATAT} (Adaptive Token Attention for Text Diffusion), which estimates token importance through a hybrid formulation combining two complementary signals:
\begin{equation}
\label{eq:hybrid}
    i^l = \lambda \cdot i^l_{\text{context}} + (1{-}\lambda) \cdot i^l_{\text{freq}}, \quad \lambda = 0.7,
\end{equation}
where $i^l_{\text{context}} = \sigmoid(\text{MLP}_\varphi(\rvh^l_{\text{GPT-2}}))$ captures contextual salience from a frozen GPT-2 anchor, and $i^l_{\text{freq}} = 1 - \log(\text{freq}(x^l)+1)/\log(\max_k\text{freq}(k)+1)$ provides a global rarity prior. The key empirical discovery is that \textit{neither component alone is sufficient}---the hybrid achieves 39.03 PPL vs.\ 40.12 (contextual-only) and 41.87 (frequency-only), demonstrating genuine signal complementarity where each component rescues the other's failure cases.

Beyond importance estimation, we introduce a \textit{balanced masking curriculum} that time-varies between two complementary regimes: at early diffusion stages (high corruption), important tokens are preserved as semantic anchors guiding reconstruction; at late stages (mostly recovered text), training signal is concentrated on important tokens to sharpen their prediction. This implements a principled curriculum~\citep{curriculum_learning,platanios2019competence} at the token level within the diffusion process.

\paragraph{Contributions.} Our contributions are fourfold:
\begin{enumerate}[leftmargin=*,itemsep=1pt,topsep=2pt]
    \item We develop a \textbf{hybrid importance estimator} combining learned contextual features (from a frozen GPT-2 anchor) with corpus frequency priors, achieving 0.82 correlation with oracle surprisal---outperforming either signal alone (0.78 and 0.71).
    \item We introduce \textbf{balanced masking}, a time-varying masking schedule that interpolates between anchor-preservation and capacity-allocation regimes, implementing curriculum learning within the diffusion process itself.
    \item We \textbf{prove theoretically} that importance-guided masking reduces sample complexity from $\gO(K^L)$ to $\gO(LK^{d+1})$ under bounded anchor assumptions, and that the balanced schedule is Lipschitz-robust to importance estimation errors.
    \item We demonstrate \textbf{state-of-the-art zero-shot perplexity} across all 7 evaluation benchmarks, training only 49M parameters (39\% of GPT-2), with particular strength on specialized domains where learned importance captures domain-specific terminology that frequency heuristics miss.
\end{enumerate}

%==============================================================================
\section{Background and Preliminaries}
\label{sec:background}
%==============================================================================

\paragraph{Notation.} Let $\gV$ denote a vocabulary of $K$ tokens, $\rvx = (x^1,\ldots,x^L) \in \gV^L$ a sequence of length $L$, and $\rvz_t \in (\gV \cup \{\texttt{[M]}\})^L$ the corrupted state at diffusion time $t \in [0,1]$.

\subsection{Absorbing-State Discrete Diffusion}

Discrete diffusion models~\citep{sohl2015deep,d3pm} define a forward process that progressively corrupts data and a reverse process that learns to denoise. For language, the \textit{absorbing-state} formulation~\citep{d3pm,mdlm} replaces tokens with a special mask token $\texttt{[M]}$. The forward marginal factorizes across positions (Eq.~\ref{eq:forward}), with noise schedule $\alpha_t$ controlling corruption severity.

\paragraph{Transition Kernel.} For $s < t$, the transition kernel from time $s$ to $t$ is $q(z_t^l \mid z_s^l) = \alpha_{t|s} \cdot \ind[z_t^l = z_s^l] + (1{-}\alpha_{t|s}) \cdot \ind[z_t^l = \texttt{[M]}]$, where $\alpha_{t|s} = \alpha_t / \alpha_s$. The absorbing property ensures that once masked, a token remains masked.

\paragraph{True Posterior.} For the reverse process, the posterior $q(z_s^l \mid z_t^l, x^l)$ is tractable. If $z_t^l = x^l$, then necessarily $z_s^l = x^l$ (since masking is irreversible). If $z_t^l = \texttt{[M]}$, then $q(z_s^l = x^l \mid z_t^l = \texttt{[M]}, x^l) = (\alpha_s - \alpha_t)/(1 - \alpha_t)$, and the remaining mass goes to $z_s^l = \texttt{[M]}$.

\paragraph{Training Objective.} DLMs are trained by minimizing the negative evidence lower bound (NELBO), which in the continuous-time limit becomes~\citep{campbell2022continuous,mdlm}:
\begin{equation}
    \gL_{\text{NELBO}}(\rvx;\theta) = \E_{t,\rvz_t} \!\left[\frac{-\dot{\alpha}_t}{1{-}\alpha_t}\!\sum_{l:z_t^l=\texttt{[M]}} \!\!-\log p_\theta(x^l|\rvz_t)\right],
    \label{eq:nelbo}
\end{equation}
where $\dot{\alpha}_t = d\alpha_t/dt$ and the sum runs only over masked positions. In practice, the ratio $-\dot{\alpha}_t/(1{-}\alpha_t)$ acts as a time-dependent weighting that emphasizes learning at intermediate noise levels.

\subsection{The Uniform Masking Limitation}

In standard DLMs, every position has identical masking probability $P(z_t^l = \texttt{[M]} \mid x^l) = 1 - \alpha_t$ regardless of position $l$ or token identity $x^l$. This treats ``quantum'' and ``the'' identically, despite the former requiring world knowledge and the latter being near-deterministic from local grammar. Information-theoretically, this is suboptimal: the model receives equal training signal from reconstructing high-entropy tokens (which require learning complex distributions) and low-entropy tokens (which can be trivially predicted from syntax alone).

ADLM~\citep{adlm} addresses this partially through frequency-based anchoring, defining anchor tokens as those with corpus frequency below a threshold. While effective (reducing PPL from 27.04 to 24.46 on LM1B relative to MDLM), frequency anchoring has inherent limitations: (1)~it cannot capture context-dependent importance variations---as noted in the ADLM work itself, the frequency proxy ``may not always align with semantic relevance''; (2)~it treats all rare tokens equally regardless of their contextual role; (3)~it is a binary decision (anchor or not) rather than a graded importance score; and (4)~it requires a dedicated anchor network (170M parameters, 12 transformer layers) in addition to the denoiser, totaling 293M trainable parameters. Our work addresses all four limitations: ATAT learns continuous, context-dependent importance scores via a lightweight 200K-parameter MLP, achieving superior results with only 49M trainable parameters---a $6\times$ reduction.

%==============================================================================
\section{Method: ATAT}
\label{sec:method}
%==============================================================================

We present ATAT in three components: the hybrid importance estimator~(\S\ref{sec:importance}), the balanced masking curriculum~(\S\ref{sec:masking}), and the full architecture and training procedure~(\S\ref{sec:arch}).

\subsection{Hybrid Importance Estimation}
\label{sec:importance}

The importance estimator $f_\varphi: \gV^L \to [0,1]^L$ assigns each token a continuous importance score by combining two complementary signals.

\paragraph{Contextual Component.} A frozen GPT-2 Small (124M parameters) serves as an anchor network, providing contextualized representations $\rvh^l_{\text{GPT-2}} = \text{GPT-2}(\rvx)^l \in \R^{768}$. A lightweight MLP transforms these into importance scores:
\begin{equation}
    i^l_{\text{context}} = \sigmoid\!\bigl(\mathbf{W}_2 \cdot \text{GELU}(\mathbf{W}_1 \cdot \text{LN}(\rvh^l_{\text{GPT-2}}) + \mathbf{b}_1) + \mathbf{b}_2\bigr),
\end{equation}
where $\mathbf{W}_1 \in \R^{256 \times 768}$, $\mathbf{W}_2 \in \R^{1 \times 256}$, and LN denotes LayerNorm. The intuition is information-theoretic: tokens that are difficult for GPT-2 to predict (high surprisal) carry more information and should receive higher importance. We supervise with oracle targets derived from GPT-2's own surprisal:
\begin{equation}
    i^{l,*} = \min\!\left(1,\; \frac{-\log p_{\text{GPT-2}}(x^l \mid x^{1:l-1})}{\tau}\right)\!, \quad \tau = 10.
    \label{eq:oracle}
\end{equation}
The temperature $\tau$ normalizes surprisal values (which range from 0 to $\log K \approx 10.8$) into $[0,1]$.

\paragraph{Frequency Component.} Corpus frequency provides a context-independent prior on token informativeness. Following Zipf's law~\citep{zipf1949human}, rare tokens tend to carry more information:
\begin{equation}
    i^l_{\text{freq}} = 1 - \frac{\log(\text{freq}(x^l)+1)}{\log(\max_k\text{freq}(k)+1)},
\end{equation}
where $\text{freq}(x^l)$ is the count of token $x^l$ in the training corpus. This assigns $i_{\text{freq}} \approx 1$ to hapax legomena and $i_{\text{freq}} \approx 0$ to the most common tokens. Unlike contextual importance, this component requires no learned parameters and provides a stable signal from the start of training, when the contextual MLP has not yet converged.

\paragraph{Hybrid Formulation.} Neither signal alone is sufficient. Frequency fails on common-but-important tokens (``the'' in emphatic constructions, frequent verbs with context-dependent significance); contextuality can be noisy early in training and may miss globally rare tokens whose importance is frequency-driven rather than context-dependent. The hybrid (Eq.~\ref{eq:hybrid}) with $\lambda = 0.7$ combines both, giving majority weight to the richer contextual signal while retaining frequency as a stabilizing prior. Theoretical justification via variance reduction is provided in Appendix~\ref{app:hybrid}.

\begin{table}[t]
\centering
\caption{Importance estimation ablation (WikiText-2 validation PPL, 50K-step ablation runs on 10\% OpenWebText subset). The hybrid consistently outperforms either component alone.}
\label{tab:hybrid-discovery}
\small
\begin{tabular}{lccc}
\toprule
\textbf{Strategy} & \textbf{PPL} & $\Delta$ vs Uniform & \textbf{Oracle Corr.} \\
\midrule
Uniform (baseline) & 42.31 & --- & --- \\
Frequency-only & 41.87 & $-$0.44 & 0.71 \\
Context-only & 40.12 & $-$2.19 & 0.78 \\
\textbf{Hybrid (ours)} & \textbf{39.03} & $\mathbf{-3.28}$ & \textbf{0.82} \\
\bottomrule
\end{tabular}
\end{table}

Table~\ref{tab:hybrid-discovery} demonstrates the complementarity: the hybrid improves by 3.28 PPL over uniform masking (7.8\% relative). Frequency captures globally rare tokens that deserve attention across all contexts; contextuality identifies locally surprising tokens whose importance depends on surrounding discourse. Qualitatively, common-but-important tokens (``the'' in ``\textit{The one who survived}'') are rescued by contextuality, while rare-but-predictable domain terms (``e.g.'' in scientific text) are stabilized by frequency.

\subsection{Curriculum-Based Balanced Masking}
\label{sec:masking}

Given importance scores $\rvi$, we design a masking schedule that varies with both token importance and diffusion timestep $t$, implementing a curriculum that shifts the model's learning focus over the course of the diffusion process.

\paragraph{Competing Objectives.} Two natural strategies for importance-aware masking are mutually exclusive:
\begin{itemize}[leftmargin=*,itemsep=1pt,topsep=2pt]
    \item \textbf{Importance-proportional}: $g_{\text{prop}}(i) = \eta + (1{-}2\eta)i$. Masks important tokens more often, providing maximal training signal on difficult predictions but destroying semantic anchors needed for coherent reconstruction.
    \item \textbf{Importance-inverse}: $g_{\text{inv}}(i) = \eta + (1{-}2\eta)(1{-}i)$. Preserves important tokens, maintaining context structure but undertraining on the hardest predictions.
\end{itemize}
Here $\eta = 0.3$ is a smoothing parameter that prevents any token from being deterministically masked or preserved, ensuring exploration across all positions.

\begin{definition}[Balanced Masking]
The balanced masking schedule interpolates between these strategies as a function of diffusion time:
\begin{equation}
\label{eq:balanced}
    g_{\text{bal}}(i,t) = t \cdot g_{\text{inv}}(i) + (1{-}t) \cdot g_{\text{prop}}(i),
\end{equation}
yielding position-dependent masking probabilities:
\begin{equation}
    q_{\text{ATAT}}(z_t^l{=}\texttt{[M]}|x^l,i^l) = (1{-}\alpha_t) \cdot g_{\text{bal}}(i^l,t).
\end{equation}
\end{definition}

\paragraph{Curriculum Interpretation.} The schedule implements a principled two-phase curriculum over the diffusion process:
\begin{itemize}[leftmargin=*,itemsep=1pt,topsep=2pt]
    \item \textbf{Heavy corruption regime} ($t \to 1$): $g_{\text{bal}} \approx g_{\text{inv}}$. Important tokens are selectively preserved as semantic anchors, providing structural scaffolding for reconstructing the many masked positions. This is analogous to providing key hints during a difficult exam---the remaining structure guides the model's predictions.
    \item \textbf{Light corruption regime} ($t \to 0$): $g_{\text{bal}} \approx g_{\text{prop}}$. Most tokens are already recovered; training signal is concentrated on important (difficult) tokens, sharpening the model's ability to predict precisely the tokens that matter most for perplexity and generation quality.
\end{itemize}


\subsection{Architecture and Training}
\label{sec:arch}

\paragraph{Anchor Network.} A frozen GPT-2 Small (124M params, 12 layers, 768 hidden dim, 12 heads) provides contextualized representations. Freezing serves three purposes: (1)~it prevents catastrophic forgetting of pretrained knowledge, (2)~it dramatically reduces trainable parameters, and (3)~it provides stable importance targets throughout training.

\paragraph{Importance Estimator.} A 2-layer MLP with 200K parameters ($768 \to 256 \to 1$, GELU activation, LayerNorm). This lightweight architecture is deliberate: the frozen GPT-2 already provides rich representations, and the MLP simply learns to extract the importance signal from them.

\paragraph{Denoiser.} A 6-layer Transformer (48M params) with the same hidden dimension (768) and attention head count (12) as GPT-2. Input representations are augmented with importance information via a learned projection:
\begin{equation}
    \tilde{\rvh}^l = \text{LN}\!\left(\rvh^l_{\text{GPT-2}} + \mathbf{W}_i[i^l;\; 1{-}i^l]\right),
\end{equation}
where $\mathbf{W}_i \in \R^{768 \times 2}$ projects the 2D importance embedding $[i^l; 1{-}i^l]$ into the representation space. The denoiser produces logits $\mathbf{o}^l$ over the vocabulary, with predictions $p_\psi(x^l|\rvz_t,\rvi) = \softmax(\mathbf{W}_{\text{out}}\mathbf{g}^l)_{x^l}$.

\paragraph{Training Objective.} The full ATAT loss combines denoising and importance supervision:
\begin{equation}
    \gL_{\text{ATAT}} = \gL_{\text{denoise}} + \gamma \gL_{\text{importance}},
    \label{eq:loss}
\end{equation}
where:
\begin{align}
    \gL_{\text{denoise}} &= \E_{t,\rvz_t}\!\left[\sum_{l:z_t^l=\texttt{[M]}} w(i^l)\bigl(-\log p_\psi(x^l|\rvz_t,\rvi)\bigr)\right], \label{eq:denoise-loss} \\
    \gL_{\text{importance}} &= \frac{1}{L}\sum_{l=1}^{L} \left(i^l_{\text{context}} - i^{l,*}\right)^2, \label{eq:imp-loss}
\end{align}
with importance-weighted cross-entropy $w(i) = 1 + \beta i$ ($\beta = 1$) upweighting difficult tokens, oracle targets $i^{l,*}$ from Eq.~\ref{eq:oracle}, and supervision weight $\gamma = 0.003$.

\paragraph{Curriculum Scheduler.} We additionally implement a macro-level training curriculum that progressively expands the importance range of tokens used for training, from easy (low-importance) to hard (full range). Details are provided in Appendix~\ref{app:masking}.

\paragraph{Parameter Summary.} Total trainable: \textbf{49M} (200K estimator + 48M denoiser), representing 39\% of GPT-2's 124M. The frozen anchor adds 124M at inference, yielding 173M total parameters.

%==============================================================================
\section{Theoretical Results}
\label{sec:theory}
%==============================================================================

We establish three theoretical properties of ATAT. Full proofs are provided in Appendix~\ref{app:theory}.

\begin{theorem}[Importance-Weighted ELBO]
\label{thm:iw-elbo}
Under the importance-guided forward process $q_{\text{ATAT}}$ with weight function $w(i) = 1 + \beta i$ ($\beta \geq 0$), the importance-weighted NELBO is a valid upper bound on the negative log-likelihood:
\begin{equation}
    -\log p_\theta(\rvx) \leq \frac{1}{\min_l w(i^l)} \cdot \gL_{\text{IW-NELBO}}(\rvx;\theta) + C,
\end{equation}
where $C$ is independent of $\theta$. Equality holds when $p_\theta = q$.
\end{theorem}

This guarantees that minimizing the ATAT objective (Eq.~\ref{eq:loss}) optimizes a valid bound on the data log-likelihood, despite the non-uniform masking.

\begin{proposition}[Sample Complexity Reduction]
\label{prop:complexity}
Under the bounded anchor assumption (Assumption~\ref{assm:anchor} in the Appendix), importance-guided masking reduces sample complexity from $\gO(K^L)$ to $\gO(LK^{d+1})$, where $d = |\gA|$ is the anchor set size.
\end{proposition}

\begin{proposition}[Robustness to Estimation Error]
\label{prop:robustness}
The balanced masking schedule (Eq.~\ref{eq:balanced}) with smoothing $\eta=0.3$ is Lipschitz continuous in importance: for estimation error $\|\hat{\rvi} - \rvi^*\|_\infty \leq \epsilon_i$, masking probability deviation is bounded by $(1{-}2\eta)\epsilon_i = 0.4\epsilon_i$.
\end{proposition}

The robustness result is practically significant: even with imperfect importance estimates (inevitable early in training), the masking schedule deviates only mildly from the oracle, preventing catastrophic training instability. The Lipschitz constant $0.4$ (reduced from the maximum possible $1.0$ by the smoothing parameter $\eta$) ensures graceful degradation.

%==============================================================================
\section{Experiments}
\label{sec:experiments}
%==============================================================================

\subsection{Experimental Setup}

\paragraph{Training Data.} We train on OpenWebText~\citep{owt}, a publicly available reproduction of the WebText corpus containing approximately 8M web documents. Sequences are tokenized using GPT-2's BPE tokenizer ($K = 50{,}257$) and packed into chunks of length $L = 1024$ by concatenating documents with $\langle \text{EOS} \rangle$ separators, following MDLM~\citep{mdlm}. No padding is used during training. Token frequencies are computed once over the full training set and cached for the frequency component of the hybrid importance estimator.

\paragraph{Training Configuration.} We train for up to 1M steps with batch size 64, using AdamW~\citep{adamw} with learning rate $3 \times 10^{-4}$, weight decay 0.01, linear warmup for 1,000 steps, and cosine decay. Gradient clipping is set to 1.0. Training uses mixed precision (FP16) on 4$\times$ NVIDIA RTX 4090 GPUs (24GB each), requiring approximately 72 hours for 1M steps. We report results at two checkpoints: 500K steps (32.8B tokens) and 1M steps (65.5B tokens). The noise schedule follows the log-linear form $\alpha_t = 1 - t$ from MDLM. Full hyperparameter tables are provided in Appendix~\ref{app:implementation}.

\paragraph{Evaluation.} We evaluate zero-shot perplexity on seven benchmarks spanning diverse domains: WikiText-2~\citep{wikitext}, LAMBADA~\citep{lambada}, Penn Treebank~\citep{ptb}, LM1B~\citep{lm1b}, AG News, PubMed, and ArXiv. Perplexity is computed using the NELBO bound with 1000 function evaluations (NFE), following~\citet{mdlm}. All evaluations use the same model checkpoint trained only on OpenWebText, without any fine-tuning or domain adaptation.

\paragraph{Baselines.} We compare against: (1)~\textbf{AR}: GPT-2 Small retrained on OpenWebText; (2)~\textbf{D3PM}~\citep{d3pm}; (3)~\textbf{SEDD}~\citep{sedd}; (4)~\textbf{MDLM}~\citep{mdlm}; (5)~\textbf{ADLM}~\citep{adlm}. All baseline numbers are from their respective papers. For fair comparison, all models use the same GPT-2 BPE tokenizer and are trained on OpenWebText.

\subsection{Main Results}

\begin{table*}[t]
\centering
\caption{Perplexity comparison across models (lower is better). $\dagger$: results from original papers. All models trained on OpenWebText. All DLM perplexities evaluated at 1000 NFE. \textbf{Train} = trainable parameters; \textbf{Tok.} = total training tokens processed.}
\label{tab:main}
\small
\resizebox{.65\textwidth}{!}{%
\begin{tabular}{llccccccc}
\toprule
\textbf{Type} & \textbf{Model} & \textbf{LM1B} & \textbf{Wiki.} & \textbf{LAM.} & \textbf{PTB} & \textbf{Params} & \textbf{Train} & \textbf{Tok.} \\
\midrule
AR & GPT-2$^\dagger$ & 21.55 & 25.75 & 51.28 & 82.05 & 124M & 124M & 65B \\
\midrule
\multirow{4}{*}{DLM} & D3PM$^\dagger$ & 47.50 & -- & -- & -- & 124M & 124M & 33B \\
& SEDD$^\dagger$ & 32.79 & 34.28 & 49.86 & 100.09 & 124M & 124M & 33B \\
& MDLM$^\dagger$ & 27.04 & 32.83 & 47.52 & 95.26 & 124M & 124M & 33B \\
& ADLM$^\dagger$ & 24.46 & 31.94 & 44.32 & 95.37 & 170M & 170M & 65B \\
\midrule
\multirow{2}{*}{Ours} & ATAT (500K) & 25.89 & 32.15 & 44.97 & 95.48 & 173M & 49M & 32.8B \\
& ATAT (1M) & \textbf{24.21} & \textbf{30.47} & \textbf{43.52} & \textbf{94.63} & 173M & 49M & 65.5B \\
\bottomrule
\end{tabular}
}
\end{table*}

Table~\ref{tab:main} presents the main perplexity results. ATAT (1M) achieves the best perplexity among all DLMs on every benchmark, and surpasses ADLM despite training \textbf{3.5$\times$ fewer parameters} (49M vs.\ 170M). Key findings:

\begin{itemize}[leftmargin=*,itemsep=1pt,topsep=2pt]
    \item \textbf{LAMBADA}: ATAT (43.52) substantially outperforms the AR baseline (51.28), demonstrating that importance-aware bidirectional context can surpass unidirectional processing on long-range dependency tasks. ATAT sets a new DLM state-of-the-art on LAMBADA, improving over ADLM's previous best of 44.32.
    \item \textbf{LM1B and WikiText}: ATAT (1M) achieves 24.21 and 30.47, improving over ADLM's 24.46 and 31.94, establishing new DLM state-of-the-art on general-domain benchmarks.
    \item \textbf{Training efficiency}: Even at 500K steps (32.8B tokens), ATAT already approaches ADLM-level performance, demonstrating efficient use of training compute. Notably, ATAT at 500K steps (49M trainable params, 32.8B tokens) matches MDLM at 1M steps (124M trainable params, 33B tokens) on LAMBADA (44.97 vs.\ 47.52), suggesting that importance-aware training extracts more signal per token than uniform masking.
    \item \textbf{Parameter efficiency}: ATAT trains only 49M parameters to surpass ADLM's 170M trainable parameters ($3.5\times$ reduction) and achieves comparable or better results than SEDD-medium (424M) and MD4-medium~\citep{shi2024simplified} (416M). The frozen GPT-2 anchor provides rich representations without gradient computation, making ATAT's per-step training cost comparable to MDLM despite the additional importance estimation. Wall-clock runtime is $\sim$72 hours for 1M steps on 4$\times$RTX 4090, with the importance MLP adding $<$3\% overhead.
\end{itemize}

\subsection{Zero-Shot Domain Transfer}

\begin{table}[t]
\centering
\caption{Zero-shot perplexity across 7 benchmarks (all models trained only on OpenWebText, evaluated at 1000 NFE). $\dagger$: from respective papers. ATAT achieves DLM SOTA on every benchmark and the best overall average.}
\label{tab:zero-shot}
\small
\resizebox{\columnwidth}{!}{%
\begin{tabular}{lcccccccc}
\toprule
\textbf{Model} & \textbf{LAM.} & \textbf{PTB} & \textbf{Wiki.} & \textbf{LM1B} & \textbf{AG} & \textbf{PubMed} & \textbf{ArXiv} & \textbf{Avg.} \\
\midrule
AR$^\dagger$ & 51.28 & 82.05 & 25.75 & 21.55 & 52.09 & 49.01 & 41.73 & 46.21 \\
MDLM$^\dagger$ & 47.52 & 95.26 & 32.83 & 27.04 & 61.15 & 41.89 & 37.37 & 49.01 \\
ADLM$^\dagger$ & 44.32 & 95.37 & 31.94 & 24.46 & 55.72 & 37.56 & 33.69 & 46.15 \\
\midrule
ATAT (1M) & \textbf{43.52} & \textbf{94.63} & \textbf{30.47} & \textbf{24.21} & \textbf{53.49} & \textbf{35.97} & \textbf{32.08} & \textbf{44.91} \\
\bottomrule
\end{tabular}
}
\end{table}

Table~\ref{tab:zero-shot} demonstrates ATAT's zero-shot generalization. ATAT achieves the lowest perplexity among all DLMs on every benchmark and \textbf{beats the AR baseline on average} (44.91 vs.\ 46.21), despite AR's advantage on WikiText and LM1B where left-to-right factorization closely matches the evaluation metric. Gains are pronounced on specialized domains:
\begin{itemize}[leftmargin=*,itemsep=1pt,topsep=2pt]
    \item \textbf{PubMed} ($-$13.04 vs.\ AR): Learned importance identifies biomedical terminology (drug names, anatomical terms) as contextually important, even when these terms have moderate frequency.
    \item \textbf{ArXiv} ($-$9.65 vs.\ AR): Mathematical notation and technical jargon are correctly prioritized by the contextual component, which recognizes their high surprisal in local context.
    \item \textbf{LAMBADA} ($-$7.76 vs.\ AR): Bidirectional context with importance-guided attention enables superior long-range prediction.
\end{itemize}

\begin{figure*}[t]
    \centering
    \includegraphics[width=1\textwidth]{figures/zero_shot_bars.pdf}
    \caption{Zero-shot perplexity across all 7 evaluation benchmarks. ATAT achieves the lowest perplexity on every domain, with the largest absolute gains on specialized domains (PubMed, ArXiv) where learned importance captures domain-specific terminology that frequency heuristics miss.}
    \label{fig:zero-shot-bars}
\end{figure*}

\subsection{Ablation Studies}
\label{sec:ablations}

\paragraph{Masking Strategy.} Table~\ref{tab:masking-ablation} isolates the contribution of balanced masking. Both importance-proportional and importance-inverse masking improve over uniform, but balanced masking achieves the best result by combining both through time-varying interpolation.

\begin{table}[t]
\centering
\caption{Masking strategy ablation (50K-step runs on 10\% OWT subset, WikiText-2 val PPL).}
\label{tab:masking-ablation}
\small
\begin{tabular}{lcc}
\toprule
\textbf{Strategy} & \textbf{PPL} $\downarrow$ & \textbf{Description} \\
\midrule
Uniform & 42.31 & Standard DLM baseline \\
Importance-prop. & 41.12 & More signal on hard tokens \\
Importance-inv. & 40.89 & Preserve semantic anchors \\
\textbf{Balanced (ours)} & \textbf{39.03} & Time-varying interpolation \\
\bottomrule
\end{tabular}
\end{table}

\paragraph{Hybrid Weight $\lambda$.} Table~\ref{tab:lambda-sensitivity} shows that ATAT is robust to the hybrid weight, with performance within 0.84 PPL across the range $[0.5, 0.9]$. The optimum at $\lambda = 0.7$ indicates contextual signals should dominate but frequency remains valuable as a stabilizer.

\begin{table}[t]
\centering
\caption{Sensitivity to hybrid weight $\lambda$ (50K-step ablation).}
\label{tab:lambda-sensitivity}
\small
\begin{tabular}{ccccc}
\toprule
$\lambda$ & \textbf{PPL} & $\Delta$ & \textbf{Corr.} & \textbf{Stability} \\
\midrule
0.0 (freq-only) & 41.87 & +2.84 & 0.71 & High \\
0.5 & 39.56 & +0.53 & 0.79 & High \\
\textbf{0.7} & \textbf{39.03} & --- & \textbf{0.82} & High \\
0.9 & 39.87 & +0.84 & 0.80 & Medium \\
1.0 (context-only) & 40.12 & +1.09 & 0.78 & Medium \\
\bottomrule
\end{tabular}
\end{table}

\paragraph{Training Dynamics.} We observe that $\gL_{\text{importance}}$ converges rapidly---within 20K steps the contextual estimator achieves $>$0.75 oracle correlation---while $\gL_{\text{denoise}}$ continues decreasing throughout training. This validates the two-loss design: the importance estimator quickly learns \textit{what matters}, allowing the denoiser to focus on \textit{how to reconstruct}. Learning curves and additional ablations (NFE trade-off, $\eta$/$\beta$/$\gamma$ sensitivity) are in Appendix~\ref{app:experiments}.

\paragraph{Component Synergy and Controls.} Incremental ablation shows that hybrid importance estimation ($-$1.42 PPL) and balanced masking ($-$2.16) are \emph{synergistic}, jointly achieving $-$3.28 (92\% of their naive sum). To control for additional parameters, replacing the frozen anchor with random importance scores yields 42.08---nearly identical to uniform MDLM (42.31)---confirming gains derive from learned importance, not added capacity. Full analysis in Appendix~\ref{app:extended-ablation}.

%==============================================================================
\section{Analysis}
\label{sec:analysis}
%==============================================================================

\paragraph{Linguistic Alignment.} We analyze learned importance scores by part-of-speech (POS) tag on WikiText-2 validation (Table~\ref{tab:pos-importance}). The scores strongly align with linguistic intuitions: content words (proper nouns: 0.78, nouns: 0.62, verbs: 0.58) receive substantially higher importance than function words (determiners: 0.21, punctuation: 0.15). This alignment emerges from training without any explicit POS supervision, validating the information-theoretic foundation.

\begin{table}[t]
\centering
\caption{Mean importance by POS category (WikiText-2 val).}
\label{tab:pos-importance}
\small
\begin{tabular}{lcc}
\toprule
\textbf{POS Category} & \textbf{Mean $i^l$} & \textbf{Linguistic Role} \\
\midrule
Proper nouns (NNP) & 0.78 & Content (entities) \\
Common nouns (NN) & 0.62 & Content (concepts) \\
Verbs (VB*) & 0.58 & Content (actions) \\
Adjectives (JJ) & 0.54 & Content (modifiers) \\
Adverbs (RB) & 0.48 & Mixed \\
Prepositions (IN) & 0.32 & Function \\
Determiners (DT) & 0.21 & Function \\
Punctuation & 0.15 & Structural \\
\bottomrule
\end{tabular}
\end{table}



\paragraph{Complementarity Analysis.} We categorize tokens into four quadrants based on their frequency and contextual importance to understand \textit{why} the hybrid outperforms either component. The hybrid's advantage stems from two specific failure modes: (1)~\textit{common-but-important} tokens ($\sim$8\% of tokens: emphatic function words, context-dependent common terms) where frequency underestimates importance, and (2)~\textit{rare-but-predictable} tokens ($\sim$5\%: domain-specific jargon that is rare globally but highly predictable in context) where contextuality overestimates importance. The hybrid correctly handles both. Full quadrant analysis is provided in Appendix~\ref{app:hybrid}.

\paragraph{Scaling Behavior.} ATAT's perplexity decreases smoothly from 32.15 (500K steps) to 30.47 (1M steps) on WikiText-2, a 5.2\% relative improvement from doubling compute. The importance estimator saturates earlier (0.80 correlation at 200K steps vs.\ 0.82 at 1M), suggesting that further gains come primarily from the denoiser learning to better exploit importance information. Analysis of loss curves and importance convergence are in Appendix~\ref{app:experiments}.

\paragraph{Use Case: Domain-Specific Generation.} ATAT's learned importance is particularly advantageous for specialized text generation. On PubMed abstracts, the model assigns high importance ($i > 0.7$) to biomedical terms such as drug names (``\textit{metformin}'', ``\textit{rituximab}'') and anatomical terms (``\textit{hippocampal}'', ``\textit{prefrontal}''), even when these tokens have moderate corpus frequency in OpenWebText. This enables coherent domain transfer: when generating biomedical text, the model preserves key clinical terminology as anchors during the early denoising phase and then refines surrounding context. On ArXiv text, mathematical notation receives elevated importance, preserving structural coherence—unlike frequency-only methods (ADLM) that treat all rare tokens equally, sometimes prioritizing peripheral words over semantically essential terms.

%==============================================================================
\section{Related Work}
\label{sec:related}
%==============================================================================

\paragraph{Discrete Diffusion for Language.} D3PM~\citep{d3pm} extended diffusion to discrete state spaces. MDLM~\citep{mdlm} simplified this to masked diffusion with improved noise schedules. SEDD~\citep{sedd} models score ratios rather than token probabilities. ADLM~\citep{adlm} introduced frequency-based anchor tokens but cannot capture context-dependent importance; its anchor network (170M) plus denoiser (123M) totals 293M trainable parameters. MD4~\citep{shi2024simplified} learns per-token vocabulary weights via REINFORCE, but these are global and static. Concurrent work explores scaling~\citep{smdm}, block-based diffusion~\citep{bd3lm}, remasking~\citep{remdm}, uniform-state diffusion~\citep{schiff2025udlm,sahoo2025duo}, discrete denoising posterior prediction~\citep{liu2025ddpd}, and discrete flow matching~\citep{gat2024discrete,campbell2024generative}. LLaDA~\citep{llada} scales masked diffusion to billion parameters. An alternative line embeds tokens into continuous space and applies Gaussian diffusion~\citep{li2022diffusion,gong2023diffuseq,gulrajani2024likelihood,dieleman2022continuous} or flow matching~\citep{zhang2024language,hu2024flow}, but suffers from a modality gap requiring ad~hoc rounding. ATAT differs from all prior work in learning continuous, context-dependent importance scores via a lightweight MLP (200K params) for token-level curriculum within the diffusion process.

\paragraph{Token Importance and Curriculum Learning.} The heterogeneous informativeness of tokens has roots in information theory~\citep{shannon1948mathematical,shannon1951prediction} and linguistics~\citep{zipf1949human}. Attention mechanisms implicitly learn token importance~\citep{clark2019does,vig2019multiscale,tenney2019bert}; information bottleneck methods~\citep{tishby2000information} and adaptive computation~\citep{graves2016adaptive,dehghani2018universal} allocate resources non-uniformly. Curriculum learning~\citep{curriculum_learning} trains models by ordering examples by difficulty, with extensions to NMT~\citep{platanios2019competence} and LMs~\citep{campos2021curriculum}. Prior work applies curriculum at the \textit{example} level; ATAT introduces curriculum at the \textit{token level within each example}.

\paragraph{Masked and Non-Autoregressive Generation.} MaskGIT~\citep{chang2022maskgit} and MAGE~\citep{li2023mage} apply iterative unmasking to image tokens, but use heuristic objectives rather than variational bounds and unmask based solely on model confidence. Non-autoregressive text methods---insertion transformers~\citep{stern2019insertion}, Levenshtein transformers~\citep{gu2019levenshtein}, masked LM generators~\citep{wang2019bert}---enable parallel generation but lack the principled corruption-denoising framework and importance-aware curriculum that ATAT provides.

%==============================================================================
\section{Conclusion}
\label{sec:conclusion}
%==============================================================================

We have presented ATAT, a framework that addresses the fundamental limitation of uniform token treatment in discrete diffusion language models. By learning hybrid importance scores combining contextual signals with frequency priors, and implementing a balanced masking curriculum that time-varies between semantic anchor preservation and capacity allocation, ATAT achieves state-of-the-art zero-shot perplexity across all 7 evaluation benchmarks while training only 49M parameters. The key insight is that importance estimation and masking curriculum are complementary: learned importance identifies \textit{which} tokens matter, while the balanced schedule determines \textit{when} to preserve versus focus on them during the diffusion process. 
% Promising future directions include: scaling laws for importance estimation with larger anchors, generation quality evaluation via MAUVE and sample entropy, dynamic timestep-dependent importance, cross-domain importance transfer, and extending importance-aware masking to multimodal discrete diffusion.

%==============================================================================
\section*{Limitations}
\label{sec:limitations}
%==============================================================================

Our evaluation focuses primarily on perplexity; generation quality metrics (MAUVE~\citep{mauve}, sample entropy) would more directly assess downstream utility but are computationally expensive to compute reliably for diffusion models. ATAT also relies on a frozen GPT-2 Small anchor, constraining importance estimation to GPT-2's pretraining knowledge---domains far from its training distribution (e.g., code, non-English text) may yield poorly calibrated surprisal signals. It remains open whether the approach scales to larger anchors or whether the importance estimator can be decoupled from a specific anchor model. Additionally, our experiments are conducted at a single model scale (49M trainable parameters); it is unclear whether the relative gains persist at significantly larger scales, though the architectural design is scale-agnostic. Finally, like all current DLMs, inference requires ${\sim}$1000 NFE, substantially slower than single-pass AR generation; our NFE ablation (Appendix~\ref{app:experiments}) shows ATAT maintains advantages at lower budgets, but the absolute cost remains a deployment barrier.

%==============================================================================
\section*{Ethics Statement}
%==============================================================================

Our work improves the efficiency and quality of non-autoregressive text generation. We do not anticipate specific negative societal impacts beyond those shared by all language models, including potential for misuse in generating misleading text. Our training data (OpenWebText) and evaluation benchmarks are publicly available and widely used in the research community.

%==============================================================================
\section*{Acknowledgments}
%==============================================================================

% Acknowledgments will be added after review.

\paragraph{Reproducibility.} We provide complete architecture specifications, hyperparameter tables, and training algorithms in the appendix (Tables~\ref{tab:arch-specs}--\ref{tab:training-config}, Algorithms~1--2). Our implementation builds on the open-source MDLM codebase, the frozen GPT-2 anchor is publicly available via HuggingFace, and all evaluation datasets are public. A detailed reproducibility checklist is included in Appendix~\ref{app:reproducibility}.

\clearpage
\bibliography{references}

\appendix
\onecolumn

\section{Theoretical Foundations}
\label{app:theory}

This appendix provides complete mathematical derivations and proofs for the theoretical results stated in the main paper.

\subsection{Absorbing-State Discrete Diffusion}
\label{app:discrete-diffusion}

We begin with a rigorous treatment of absorbing-state discrete diffusion, establishing the framework upon which ATAT builds. Continuous diffusion~\citep{ddpm,score_sde} operates in $\R^d$ by adding Gaussian noise: $q(\rvz_t|\rvx) = \gN(\sqrt{\bar\alpha_t}\rvx, (1{-}\bar\alpha_t)\mathbf{I})$, which is natural for images where pixel values are continuous. Language tokens inhabit a discrete space $\gV^L$ where Gaussian perturbation is ill-defined. Discrete diffusion resolves this by replacing additive noise with \textbf{absorbing-state masking}: tokens are stochastically replaced with a special $\texttt{[M]}$ symbol.

\subsubsection{Forward Process}

Given clean data $\rvx = (x^1, \ldots, x^L) \in \gV^L$, the forward process defines a Markov chain that progressively corrupts tokens to an absorbing mask state. Let $\alpha_t: [0,1] \to [0,1]$ be a monotonically decreasing noise schedule with $\alpha_0 = 1$ (clean) and $\alpha_1 = 0$ (fully masked).

\begin{definition}[Forward Marginal]
The marginal distribution at time $t$ factorizes across positions:
\begin{equation}
    q(\rvz_t \mid \rvx) = \prod_{l=1}^{L} q(z^l_t \mid x^l), \quad q(z^l_t \mid x^l) = \alpha_t \cdot \delta_{z^l_t, x^l} + (1{-}\alpha_t) \cdot \delta_{z^l_t, \texttt{[M]}}
\end{equation}
\end{definition}

Each token independently survives with probability $\alpha_t$ or is replaced by $\texttt{[M]}$ with probability $1{-}\alpha_t$. The factorization across positions is a key simplification: tokens are corrupted independently, though the reverse (denoising) process models their joint distribution.

\begin{lemma}[Transition Kernel]
\label{lem:transition}
For $s < t$, the transition kernel is:
\begin{equation}
    q(\rvz_t \mid \rvz_s) = \prod_{l=1}^{L} \left[\alpha_{t|s} \cdot \delta_{z^l_t, z^l_s} + (1{-}\alpha_{t|s}) \cdot \delta_{z^l_t, \texttt{[M]}}\right]
\end{equation}
where $\alpha_{t|s} := \alpha_t / \alpha_s$.
\end{lemma}

\begin{proof}
By the law of total probability:
\begin{align}
    q(z^l_t = x^l \mid x^l) &= q(z^l_t = x^l \mid z^l_s = x^l) \cdot q(z^l_s = x^l \mid x^l) + q(z^l_t = x^l \mid z^l_s = \texttt{[M]}) \cdot q(z^l_s = \texttt{[M]} \mid x^l)
\end{align}
Since $\texttt{[M]}$ is absorbing, $q(z^l_t = x^l \mid z^l_s = \texttt{[M]}) = 0$. Thus:
\begin{equation}
    \alpha_t = q(z^l_t = x^l \mid z^l_s = x^l) \cdot \alpha_s \implies q(z^l_t = x^l \mid z^l_s = x^l) = \frac{\alpha_t}{\alpha_s} = \alpha_{t|s}
\end{equation}
This confirms that the conditional survival probability from time $s$ to $t$ is $\alpha_{t|s}$, and the conditional masking probability is $1 - \alpha_{t|s}$.
\end{proof}

\subsubsection{True Posterior}

The true posterior $q(z^l_s \mid z^l_t, x^l)$ is required for the KL divergence computation in the ELBO.

\begin{lemma}[Posterior Distribution]
\label{lem:posterior}
For $s < t$, the posterior satisfies:
\begin{equation}
    q(z^l_s \mid z^l_t, x^l) = \begin{cases}
        \delta_{z^l_s, x^l} & \text{if } z^l_t = x^l \\[0.5em]
        \frac{\alpha_s - \alpha_t}{1 - \alpha_t} \cdot \delta_{z^l_s, x^l} + \frac{1 - \alpha_s}{1 - \alpha_t} \cdot \delta_{z^l_s, \texttt{[M]}} & \text{if } z^l_t = \texttt{[M]}
    \end{cases}
\end{equation}
\end{lemma}

\begin{proof}
By Bayes' theorem:
\begin{equation}
    q(z^l_s \mid z^l_t, x^l) = \frac{q(z^l_t \mid z^l_s) \cdot q(z^l_s \mid x^l)}{q(z^l_t \mid x^l)}
\end{equation}

\textbf{Case 1:} $z^l_t = x^l$. Since $\texttt{[M]}$ is absorbing, we must have $z^l_s = x^l$, giving $q(z^l_s = x^l \mid z^l_t = x^l, x^l) = 1$.

\textbf{Case 2:} $z^l_t = \texttt{[M]}$. Then:
\begin{align}
    q(z^l_s = x^l \mid z^l_t = \texttt{[M]}, x^l) &= \frac{q(z^l_t = \texttt{[M]} \mid z^l_s = x^l) \cdot q(z^l_s = x^l \mid x^l)}{q(z^l_t = \texttt{[M]} \mid x^l)} \\
    &= \frac{(1 - \alpha_{t|s}) \cdot \alpha_s}{1 - \alpha_t} = \frac{\alpha_s - \alpha_t}{1 - \alpha_t}
\end{align}
The remaining probability mass goes to $z^l_s = \texttt{[M]}$:
\begin{equation}
    q(z^l_s = \texttt{[M]} \mid z^l_t = \texttt{[M]}, x^l) = 1 - \frac{\alpha_s - \alpha_t}{1 - \alpha_t} = \frac{1 - \alpha_s}{1 - \alpha_t} \qedhere
\end{equation}
\end{proof}

\subsubsection{Continuous-Time NELBO}

In the continuous-time limit ($T \to \infty$ discrete steps), the NELBO becomes~\citep{campbell2022continuous,mdlm}:
\begin{equation}
    \gL_{\text{NELBO}}(\rvx;\theta) = \int_0^1 \frac{-\dot\alpha_t}{1-\alpha_t} \E_{\rvz_t \sim q(\cdot|\rvx)}\left[\sum_{l:z^l_t=\texttt{[M]}} -\log p_\theta(x^l | \rvz_t)\right] dt
\end{equation}

The weighting $-\dot\alpha_t/(1-\alpha_t)$ arises from the rate at which new tokens are masked at time $t$. Under a log-linear schedule $\alpha_t = 1-t$ (used in MDLM), this simplifies to $1/(1-\alpha_t) = 1/t$, giving higher weight to late timesteps where few tokens remain masked and reconstruction is most informative.

\subsection{Contrast with Other Diffusion Paradigms}
\label{app:contrast}

ATAT builds on the absorbing-state framework above. Two alternative paradigms exist, each with distinct trade-offs:

\paragraph{Continuous Diffusion for Language.} An alternative line of work~\citep{li2022diffusion,gong2023diffuseq,gulrajani2024likelihood,dieleman2022continuous} embeds discrete tokens into a continuous space $\R^d$ (e.g., via pretrained word embeddings) and applies standard Gaussian diffusion. At generation time, the continuous output is \emph{rounded} back to the nearest token embedding. While this enables the use of well-studied continuous diffusion machinery (score matching~\citep{score_sde}, DDPM~\citep{ddpm}), the approach introduces a \textbf{modality gap}: the rounding step is non-differentiable and lossy, and small perturbations in continuous space can map to entirely different tokens. Recent flow-matching approaches~\citep{zhang2024language,hu2024flow} partially mitigate this but still require ad-hoc embedding/rounding procedures. ATAT avoids this issue entirely by operating natively in discrete token space.

\paragraph{Score-Entropy Discrete Diffusion (SEDD).} SEDD~\citep{sedd} proposes an alternative to the NELBO by modeling concrete score ratios $s_\theta(\rvz_t, t)_k^l = p_\theta(x^l = k \mid \rvz_t) / p_\theta(z_t^l \mid \rvz_t)$ and training with a score-entropy loss. SEDD achieves strong results on LM1B (32.79 PPL) but is less directly comparable to the absorbing-state framework because it (i)~uses a different noise process (uniform rather than absorbing), (ii)~requires score-ratio parameterization rather than direct token prediction, and (iii)~does not naturally support position-dependent masking, making it difficult to incorporate importance-aware curricula. ATAT's absorbing-state formulation cleanly supports position-dependent masking rates through Eq.~\ref{eq:balanced} in the main paper.

\subsection{Proof of Theorem~\ref{thm:iw-elbo}: Validity of Importance-Weighted ELBO}
\label{app:elbo-proof}

\begin{theorem}[Importance-Weighted ELBO --- Restated]
Let $\rvi = (i^1, \ldots, i^L)$ be importance scores with $i^l \in [0,1]$. The importance-weighted objective:
\begin{equation}
    \gL_{\text{IW}}(\rvx; \theta) = \E_{t, \rvz_t}\left[\sum_{l: z^l_t = \texttt{[M]}} w(i^l) \cdot \left(-\log p_\theta(x^l \mid \rvz_t)\right)\right]
\end{equation}
with $w(i^l) = 1 + \beta \cdot i^l$ for $\beta \geq 0$, satisfies:
\begin{equation}
    -\log p_\theta(\rvx) \leq \frac{1}{\min_l w(i^l)} \cdot \gL_{\text{IW}}(\rvx; \theta) + C
\end{equation}
where $C$ is a constant independent of $\theta$.
\end{theorem}

\begin{proof}
We verify three conditions ensuring the ELBO remains valid under adaptive masking.

\textbf{(1) Well-definedness.} The modified forward process with position-dependent masking rates $\gamma^l_t = \gamma_t \cdot g_{\text{bal}}(i^l, t)$ remains a valid probability distribution. Since $g_{\text{bal}}(i^l, t) \in [\eta, 1{-}\eta]$ for $\eta = 0.3$ and $\gamma_t \in [0,1]$, we have $\gamma^l_t \in [0, 1{-}\eta] \subset [0,1]$, ensuring valid masking probabilities.

\textbf{(2) Terminal condition.} At $t = 1$, we require all tokens to be masked with high probability. With $\gamma_1 = 1$:
\begin{equation}
    q(z^l_1 = \texttt{[M]} \mid x^l) = \gamma^l_1 = g_{\text{bal}}(i^l, 1) \geq \eta = 0.3
\end{equation}
For $\eta < 0.5$, we can rescale to ensure full masking at $t=1$ by setting $g_{\text{bal}}(i, 1) = 1$ as a boundary condition.

\textbf{(3) ELBO derivation.} Starting from the standard NELBO:
\begin{align}
    -\log p_\theta(\rvx) &\leq \E_q\left[-\log p_\theta(\rvx \mid \rvz_0)\right] + \sum_{t=1}^T \E_q\left[D_{\text{KL}}(q(\rvz_{t-1} \mid \rvz_t, \rvx) \| p_\theta(\rvz_{t-1} \mid \rvz_t))\right]
\end{align}

Introducing importance weights $w(i^l)$ with normalization $\bar{w} = \frac{1}{L}\sum_l w(i^l)$:
\begin{equation}
    \gL_{\text{IW}} = \bar{w} \cdot \E_q\left[\sum_l \frac{w(i^l)}{\bar{w}} \cdot \ell(x^l, \rvz_t)\right]
\end{equation}

Since $w(i^l) \geq w_{\min} = 1$ (for $\beta \geq 0$), we have:
\begin{equation}
    \gL_{\text{standard}} = \E_q\left[\sum_l \ell(x^l, \rvz_t)\right] \leq \frac{1}{w_{\min}} \cdot \gL_{\text{IW}}
\end{equation}

The ELBO bound follows by combining with the standard derivation. The constant $C$ absorbs the KL divergence between the modified and standard forward processes, which is bounded since both share the same absorbing state.
\end{proof}

\begin{corollary}[Tightness]
The bound is tightest when importance weights are uniform ($\beta = 0$), recovering the standard NELBO. As $\beta$ increases, the bound loosens but training focuses more on important tokens, improving practical performance despite a weaker theoretical guarantee.
\end{corollary}

\subsection{Proof of Proposition~\ref{prop:complexity}: Sample Complexity Reduction}
\label{app:complexity-proof}

\begin{assumption}[Bounded Anchor Structure]
\label{assm:anchor}
There exists a subset $\gA \subset \{1, \ldots, L\}$ of ``anchor'' positions with $|\gA| = d \ll L$ such that for all non-anchor positions $l \notin \gA$:
\begin{equation}
    p(x^l \mid \rvx^{\setminus l}) = p(x^l \mid \rvx^{\gA})
\end{equation}
i.e., non-anchor tokens depend only on the anchor tokens, not on other non-anchor tokens.
\end{assumption}

This assumption is motivated by the empirical observation that natural language has a relatively small set of ``pivotal'' tokens per sentence that determine the meaning, while remaining tokens are largely predictable from these pivots.

\begin{proposition}[Sample Complexity under Anchored Structure --- Restated]
Let $\rvx \in \gV^L$ with $|\gV| = K$. Under uniform masking, the sample complexity for $\epsilon$-accurate MLE is $\gO(K^L/\epsilon^2)$. Under importance-guided masking with bounded anchor set $|\gA| = d \ll L$, the complexity reduces to $\gO(L \cdot K^{d+1}/\epsilon^2)$.
\end{proposition}

\begin{proof}
We analyze the parameter count and apply standard MLE sample complexity bounds.

\textbf{Uniform masking.} Each conditional $p(x^l \mid \rvx^{\setminus l})$ requires specifying a distribution over $K$ outcomes for each of $K^{L-1}$ contexts, yielding $\gO(K^L)$ parameters per position and $\gO(L \cdot K^L)$ total.

\textbf{Importance-guided masking.} Under Assumption~\ref{assm:anchor}, important tokens form an anchor set $\gA$ with $|\gA| = d$. The joint factorizes as:
\begin{equation}
    p(\rvx) = p(\rvx^{\gA}) \cdot \prod_{l \notin \gA} p(x^l \mid \rvx^{\gA})
\end{equation}

Parameter counts:
\begin{itemize}
    \item Anchor distribution $p(\rvx^{\gA})$: $K^d$ parameters
    \item Each conditional $p(x^l \mid \rvx^{\gA})$: $K^{d+1}$ parameters
    \item Total: $K^d + (L{-}d) \cdot K^{d+1} = \gO(L \cdot K^{d+1})$
\end{itemize}

By standard MLE theory, sample complexity scales linearly with parameter count divided by $\epsilon^2$, giving the stated bounds.
\end{proof}

\begin{remark}
For $L = 1024$, $K = 50{,}257$, and $d = 10$: uniform masking requires $\gO(50257^{1024})$ samples, while anchored masking requires $\gO(1024 \times 50257^{11}) \approx 10^{55}$---still astronomical but exponentially better. In practice, neural networks with weight sharing dramatically reduce the effective parameter count, but the \textit{relative} improvement from importance-guided masking persists.
\end{remark}

\subsection{Proof of Proposition~\ref{prop:robustness}: Robustness to Estimation Error}
\label{app:robustness-proof}

\begin{proposition}[Lipschitz Continuity of Balanced Masking --- Restated]
Let $\hat{\rvi}$ and $\rvi^*$ be estimated and oracle importance with $\|\hat{\rvi} - \rvi^*\|_\infty \leq \epsilon_i$. Then:
\begin{equation}
    \left|g_{\text{bal}}(\hat{i}^l, t) - g_{\text{bal}}(i^{l,*}, t)\right| \leq (1 - 2\eta) \cdot \epsilon_i
\end{equation}
\end{proposition}

\begin{proof}
The balanced masking function is:
\begin{equation}
    g_{\text{bal}}(i, t) = t \cdot g_{\text{inv}}(i) + (1{-}t) \cdot g_{\text{prop}}(i)
\end{equation}
with $g_{\text{inv}}(i) = \eta + (1{-}2\eta)(1{-}i)$ and $g_{\text{prop}}(i) = \eta + (1{-}2\eta)i$.

Computing the partial derivative:
\begin{align}
    \frac{\partial g_{\text{bal}}}{\partial i} &= t \cdot \frac{\partial g_{\text{inv}}}{\partial i} + (1{-}t) \cdot \frac{\partial g_{\text{prop}}}{\partial i} \\
    &= t \cdot (-(1{-}2\eta)) + (1{-}t) \cdot (1{-}2\eta) \\
    &= (1{-}2\eta)(1 - 2t)
\end{align}

The maximum absolute value of this derivative is:
\begin{equation}
    \max_{t \in [0,1]} |(1{-}2\eta)(1-2t)| = |1{-}2\eta| \cdot \max_{t \in [0,1]} |1-2t| = (1{-}2\eta) \cdot 1
\end{equation}
achieved at $t \in \{0, 1\}$. By the mean value theorem:
\begin{equation}
    |g_{\text{bal}}(\hat{i}^l, t) - g_{\text{bal}}(i^{l,*}, t)| \leq (1{-}2\eta) \cdot |\hat{i}^l - i^{l,*}| \leq (1{-}2\eta) \cdot \epsilon_i
\end{equation}

For $\eta = 0.3$: Lipschitz constant $= 1 - 0.6 = 0.4$, so masking probabilities change by at most $0.4 \epsilon_i$.
\end{proof}


The smoothing parameter $\eta = 0.3$ serves dual purposes: (1) it ensures exploration by preventing deterministic masking/preservation ($g_{\text{bal}} \in [0.3, 0.7]$), and (2) it reduces sensitivity to importance estimation errors by a factor of $2.5\times$ compared to the unsmoothed case ($\eta = 0$). This is critical during early training when the importance estimator has not yet converged.


\subsection{Additional Theoretical Results}

\begin{lemma}[Conservation of Expected Masking Rate]
\label{lem:conservation}
Under balanced masking with $\E[i^l] = \mu$:
\begin{equation}
    \E_l[g_{\text{bal}}(i^l, t)] = \eta + (1{-}2\eta)\left[(1{-}\mu)t + \mu(1{-}t)\right]
\end{equation}
When $\mu = 0.5$ (symmetric importance distribution), $\E_l[g_{\text{bal}}(i^l, t)] = 0.5$ for all $t$.
\end{lemma}

\begin{proof}
\begin{align}
    \E_l[g_{\text{bal}}(i^l, t)] &= t \cdot \E_l[g_{\text{inv}}(i^l)] + (1{-}t) \cdot \E_l[g_{\text{prop}}(i^l)] \\
    &= t \cdot [\eta + (1{-}2\eta)(1{-}\mu)] + (1{-}t) \cdot [\eta + (1{-}2\eta)\mu] \\
    &= \eta + (1{-}2\eta)[(1{-}\mu)t + \mu(1{-}t)]
\end{align}
Substituting $\mu = 0.5$: $(1{-}0.5)t + 0.5(1{-}t) = 0.5t + 0.5 - 0.5t = 0.5$, so $\E_l[g_{\text{bal}}] = \eta + (1{-}2\eta) \cdot 0.5 = 0.5$.
\end{proof}

This conservation property means balanced masking redistributes masking probability across tokens without changing the total expected number of masked tokens, ensuring compatibility with the standard NELBO training framework.

\begin{lemma}[Monotonicity of Curriculum Effect]
\label{lem:monotonicity}
For a token with importance $i > 0.5$:
\begin{itemize}
    \item $g_{\text{bal}}(i, t)$ is strictly decreasing in $t$ (less masking at higher corruption)
    \item At $t=1$: $g_{\text{bal}}(i, 1) = g_{\text{inv}}(i) < 0.5$ (preserved as anchor)
    \item At $t=0$: $g_{\text{bal}}(i, 0) = g_{\text{prop}}(i) > 0.5$ (focused for training)
\end{itemize}
For tokens with $i < 0.5$, the behavior reverses.
\end{lemma}

\begin{proof}
\begin{equation}
    \frac{\partial g_{\text{bal}}}{\partial t} = g_{\text{inv}}(i) - g_{\text{prop}}(i) = (1{-}2\eta)(1{-}2i)
\end{equation}
For $i > 0.5$ and $\eta < 0.5$: $\partial g_{\text{bal}}/\partial t = (1{-}2\eta)(1{-}2i) < 0$, confirming the decreasing property.
\end{proof}

\section{Hybrid Importance Estimation: Theoretical Analysis}
\label{app:hybrid}

This section provides the theoretical foundation for our hybrid importance estimator, analyzing why the combination of frequency-based and contextual signals outperforms either alone.

\subsection{Information-Theoretic Foundation}

\begin{definition}[Token Importance]
The \emph{oracle importance} of token $x^l$ in context $\rvx^{\setminus l}$ is:
\begin{equation}
    i^{l,*} = \frac{H[X^l \mid \rvx^{\setminus l}]}{\log K}
\end{equation}
where $H[\cdot]$ denotes entropy and $K = |\gV|$.
\end{definition}

This normalizes importance to $[0,1]$: maximally uncertain tokens have $i^* = 1$; deterministic tokens have $i^* = 0$. In practice, we approximate this with GPT-2 surprisal, which uses left context only ($\rvx^{<l}$ instead of $\rvx^{\setminus l}$).

\subsection{Frequency-Based Approximation}

\begin{proposition}[Frequency as Importance Proxy]
\label{prop:freq-proxy}
Under a unigram model $p_{\text{unigram}}(x) = \text{freq}(x)/N$, the frequency-based importance:
\begin{equation}
    i^l_{\text{freq}} = 1 - \frac{\log(\text{freq}(x^l) + 1)}{\log(\text{max\_freq} + 1)}
\end{equation}
approximates $i^{l,*}$ with error bounded by the mutual information $I(X^l; \rvx^{\setminus l})$.
\end{proposition}

\begin{proof}
Under the unigram assumption:
\begin{equation}
    H[X^l \mid \rvx^{\setminus l}] = H[X^l] = -\sum_x p(x) \log p(x)
\end{equation}

For token $x^l$ with frequency $f = \text{freq}(x^l)$:
\begin{equation}
    -\log p(x^l) = -\log(f/N) = \log N - \log f \approx \log(\text{max\_freq}) - \log f
\end{equation}

Normalizing: $i_{\text{freq}} \approx 1 - \log f / \log(\text{max\_freq})$, which matches our formula.

The approximation error is:
\begin{equation}
    |i^{l,*} - i_{\text{freq}}| \leq \frac{I(X^l; \rvx^{\setminus l})}{\log K}
\end{equation}
capturing the context-dependence not modeled by the unigram assumption. This error is large for context-dependent tokens (where knowing the context significantly changes the prediction) and small for context-independent tokens.
\end{proof}

\subsection{Contextual Importance via Surprisal}

\begin{proposition}[Surprisal as Contextual Importance]
The GPT-2 surprisal-based importance:
\begin{equation}
    i^l_{\text{context}} = \sigmoid\left(\text{MLP}_\varphi\left(\rvh^l_{\text{GPT-2}}\right)\right)
\end{equation}
trained with MSE loss against oracle surprisal $-\log p_{\text{GPT-2}}(x^l \mid \rvx^{<l})$, converges to an estimator of $H[X^l \mid \rvx^{<l}]$.
\end{proposition}

\begin{proof}
The training objective minimizes:
\begin{equation}
    \gL = \E\left[\left(i^l_{\text{context}} - \frac{-\log p_{\text{GPT-2}}(x^l \mid \rvx^{<l})}{\tau}\right)^2\right]
\end{equation}

At convergence, the optimal predictor satisfies:
\begin{equation}
    i^l_{\text{context}} = \E\left[\frac{-\log p(x^l \mid \rvx^{<l})}{\tau} \;\middle|\; \rvh^l_{\text{GPT-2}}\right] = \frac{H[X^l \mid \rvx^{<l}]}{\tau}
\end{equation}
assuming the MLP has sufficient capacity and $\rvh^l_{\text{GPT-2}}$ is a sufficient statistic for the left context.
\end{proof}

\subsection{Hybrid Combination: Variance Reduction}

\begin{theorem}[Variance Reduction via Hybrid Estimation]
\label{thm:hybrid-variance}
Let $\hat{i}_{\text{freq}}$ and $\hat{i}_{\text{context}}$ be unbiased estimators of $i^*$ with variances $\sigma^2_{\text{freq}}$ and $\sigma^2_{\text{context}}$ and covariance $\rho \sigma_{\text{freq}} \sigma_{\text{context}}$. The hybrid estimator:
\begin{equation}
    \hat{i}_{\text{hybrid}} = \lambda \cdot \hat{i}_{\text{context}} + (1{-}\lambda) \cdot \hat{i}_{\text{freq}}
\end{equation}
has minimum variance at:
\begin{equation}
    \lambda^* = \frac{\sigma^2_{\text{freq}} - \rho \sigma_{\text{freq}} \sigma_{\text{context}}}{\sigma^2_{\text{freq}} + \sigma^2_{\text{context}} - 2\rho \sigma_{\text{freq}} \sigma_{\text{context}}}
\end{equation}
\end{theorem}

\begin{proof}
The variance of the hybrid estimator is:
\begin{align}
    \text{Var}[\hat{i}_{\text{hybrid}}] &= \lambda^2 \sigma^2_{\text{context}} + (1{-}\lambda)^2 \sigma^2_{\text{freq}} + 2\lambda(1{-}\lambda)\rho \sigma_{\text{freq}} \sigma_{\text{context}}
\end{align}
Taking the derivative with respect to $\lambda$ and setting to zero:
\begin{align}
    \frac{d}{d\lambda}\text{Var}[\hat{i}_{\text{hybrid}}] &= 2\lambda\sigma^2_{\text{context}} - 2(1{-}\lambda)\sigma^2_{\text{freq}} + 2(1{-}2\lambda)\rho\sigma_{\text{freq}}\sigma_{\text{context}} = 0
\end{align}
Solving for $\lambda$ yields the stated $\lambda^*$.
\end{proof}

\begin{corollary}
When frequency and contextuality errors are uncorrelated ($\rho = 0$) and have equal variance ($\sigma_{\text{freq}} = \sigma_{\text{context}} = \sigma$), then $\lambda^* = 0.5$ and $\text{Var}[\hat{i}_{\text{hybrid}}] = \sigma^2/2$---a $2\times$ variance reduction over either estimator alone.
\end{corollary}

\begin{remark}
Our empirical $\lambda = 0.7$ (favoring contextuality) suggests that after training convergence, $\sigma^2_{\text{context}} < \sigma^2_{\text{freq}}$---the learned contextual estimator has lower variance than the fixed frequency prior. This makes sense: the contextual component adapts to the specific token-context pair, while frequency provides only a global prior.
\end{remark}

\subsection{Complementarity Analysis}

We categorize tokens into four quadrants to understand when each component contributes:

\begin{table}[h]
\centering
\caption{Four-quadrant analysis of hybrid importance components.}
\label{tab:quadrant}
\begin{tabular}{p{3.5cm}p{3.5cm}p{3.5cm}p{3.5cm}}
\toprule
& \textbf{High Freq. Importance} & \textbf{Low Freq. Importance} \\
\midrule
\textbf{High Context Importance} & Rare, surprising tokens (``quantum'' in news): Both signals agree. $\sim$35\% of tokens. & Common but surprising (``the'' in emphasis): Frequency underestimates, context rescues. $\sim$8\%. \\
\midrule
\textbf{Low Context Importance} & Rare but predictable (jargon in domain text): Context overestimates, frequency stabilizes. $\sim$5\%. & Common, predictable (``the'' normally): Both signals agree. $\sim$52\%. \\
\bottomrule
\end{tabular}
\end{table}

The hybrid's 3.28 PPL advantage over uniform masking comes disproportionately from the $\sim$13\% of tokens in the off-diagonal quadrants, where the two signals provide complementary information.

\section{Complete Algorithm Specifications}
\label{app:algorithms}

\subsection{Training Algorithm}

\begin{algorithm}[H]
\caption{ATAT Training Procedure}
\label{alg:atat-train}
\begin{algorithmic}[1]
\REQUIRE Dataset $\gD$, frozen anchor $\gA$ (GPT-2), learning rate $\alpha_{\text{lr}}$, total steps $S = 1{,}000{,}000$
\REQUIRE Hyperparameters: $\lambda=0.7$, $\beta=1.0$, $\gamma=0.003$, $\eta=0.3$
\STATE Initialize importance estimator $f_\varphi$ (200K params)
\STATE Initialize denoiser $\gD_\psi$ from GPT-2 layers 1--6 (48M params)
\STATE Precompute token frequencies $\{\text{freq}(v)\}_{v \in \gV}$ from $\gD$
\FOR{$s = 1$ to $S$}
    \STATE Sample batch $\{\rvx_b\}_{b=1}^{B} \sim \gD$ \hfill ($B = 64$, $L = 1024$)
    \STATE $\rvh_b = \gA(\rvx_b)$ \hfill \textit{// Anchor representations (frozen)}
    \STATE $\rvi^{\text{context}}_b = \sigmoid(\text{MLP}_\varphi(\rvh_b))$ \hfill \textit{// Contextual importance}
    \STATE $i^{l,\text{freq}}_b = 1 - \log(\text{freq}(x^l_b){+}1)/\log(\text{max\_freq}{+}1)$ \hfill \textit{// Frequency importance}
    \STATE $\rvi_b = \lambda \cdot \rvi^{\text{context}}_b + (1{-}\lambda) \cdot \rvi^{\text{freq}}_b$ \hfill \textit{// Hybrid combination}
    \STATE $t \sim \text{Uniform}[0, 1]$
    \STATE $g^l_b = t \cdot g_{\text{inv}}(i^l_b) + (1{-}t) \cdot g_{\text{prop}}(i^l_b)$ \hfill \textit{// Balanced schedule}
    \STATE $m^l_b \sim \text{Bernoulli}((1{-}\alpha_t) \cdot g^l_b)$ \hfill \textit{// Adaptive masking}
    \STATE $\rvz_t = \rvx_b \odot (1{-}\rvm_b) + \texttt{[M]} \odot \rvm_b$
    \STATE $\hat{\rvx}_b = \gD_\psi(\gA(\rvz_t), \rvi_b)$ \hfill \textit{// Denoising prediction}
    \STATE $\gL_{\text{denoise}} = \sum_{l: m^l_b=1} (1 + \beta \cdot i^l_b) \cdot \text{CE}(\hat{x}^l_b, x^l_b)$
    \STATE $\rvi^* = \min(1, -\log p_{\gA}(\rvx_b)/\tau)$ \hfill \textit{// Oracle targets}
    \STATE $\gL_{\text{imp}} = \|\rvi^{\text{context}}_b - \rvi^*\|_2^2 / L$
    \STATE $(\varphi, \psi) \leftarrow (\varphi, \psi) - \alpha_{\text{lr}} \nabla(\gL_{\text{denoise}} + \gamma \cdot \gL_{\text{imp}})$
\ENDFOR
\RETURN $(\varphi, \psi)$
\end{algorithmic}
\end{algorithm}

\subsection{Inference Algorithm}

\begin{algorithm}[H]
\caption{ATAT Generation with Uncertainty-Guided Sampling}
\label{alg:atat-generate}
\begin{algorithmic}[1]
\REQUIRE Trained model $(\gA, f_\varphi, \gD_\psi)$, sequence length $L$, denoising steps $T$
\STATE Initialize $\rvz_T = \texttt{[M]}^L$ \hfill \textit{// Start fully masked}
\FOR{$t = T, T{-}1, \ldots, 1$}
    \STATE $\rvh = \gA(\rvz_t)$ \hfill \textit{// Anchor representations}
    \STATE $\rvi = \lambda \cdot f_\varphi(\rvh) + (1{-}\lambda) \cdot \rvi^{\text{freq}}$ \hfill \textit{// Hybrid importance}
    \STATE $\kappa = 1/(t/T + 0.1)$ \hfill \textit{// Sharpening coefficient}
    \STATE $\rvi \leftarrow \rvi^{\kappa}$ \hfill \textit{// Sharpen importance as generation progresses}
    \STATE $\hat{\rvx}, \mathbf{o} = \gD_\psi(\rvh, \rvi)$ \hfill \textit{// Predict tokens and logits}
    \STATE $u^l = H[\softmax(\mathbf{o}^l)]$ for all masked $l$ \hfill \textit{// Predictive entropy}
    \STATE $s^l = u^l \cdot i^l$ for all masked $l$ \hfill \textit{// Priority = uncertainty $\times$ importance}
    \STATE $n = \max(1, \lfloor |\{l: z^l_t = \texttt{[M]}\}| / T \rfloor)$ \hfill \textit{// Tokens to unmask}
    \STATE $\mathcal{K} = \text{top-}n(s^l)$ over masked positions \hfill \textit{// Select highest priority}
    \FOR{$l \in \mathcal{K}$}
        \STATE $z^l_{t-1} = \argmax_k \softmax(\mathbf{o}^l)_k$ \hfill \textit{// Greedy decoding}
    \ENDFOR
    \STATE Keep remaining masked positions unchanged: $z^l_{t-1} = \texttt{[M]}$
\ENDFOR
\RETURN $\rvz_0$
\end{algorithmic}
\end{algorithm}

\subsection{Inference Design Choices}

The uncertainty-guided sampling order differs from alternatives considered:
\begin{itemize}
    \item \textbf{Random order}: Ignores both model confidence and token importance, leading to suboptimal denoising trajectories.
    \item \textbf{Confidence-first} (unmask most confident first): May denoise trivial tokens early, wasting steps on function words while important tokens remain masked.
    \item \textbf{Importance-only} (unmask most important first): Ignores whether the model is actually confident enough to make correct predictions.
    \item \textbf{Uncertainty $\times$ importance} (ours): Prioritizes positions that are both important and uncertain, focusing denoising effort where it matters most and where the model has actionable information.
\end{itemize}

The sharpening coefficient $\kappa = 1/(t/T + 0.1)$ increases as generation progresses, making importance distinctions sharper when fewer tokens remain masked and precise allocation matters more.

\section{Masking Strategy Analysis}
\label{app:masking}

\subsection{Strategy Definitions}

Given importance $\rvi \in [0,1]^L$ and smoothing parameter $\eta = 0.3$:

\paragraph{Importance-Proportional.} High importance $\to$ high masking probability:
\begin{equation}
    g_{\text{prop}}(i) = \eta + (1{-}2\eta) \cdot i
\end{equation}
Range: $[\eta, 1{-}\eta] = [0.3, 0.7]$. Maximally important tokens ($i=1$) are masked with probability $0.7$; minimally important tokens ($i=0$) with probability $0.3$.

\paragraph{Importance-Inverse.} High importance $\to$ low masking probability:
\begin{equation}
    g_{\text{inv}}(i) = \eta + (1{-}2\eta) \cdot (1{-}i)
\end{equation}
Range: $[\eta, 1{-}\eta] = [0.3, 0.7]$. Maximally important tokens ($i=1$) are masked with probability $0.3$; minimally important tokens ($i=0$) with probability $0.7$.

\paragraph{Balanced.} Time-varying interpolation:
\begin{equation}
    g_{\text{bal}}(i, t) = t \cdot g_{\text{inv}}(i) + (1{-}t) \cdot g_{\text{prop}}(i)
\end{equation}

\subsection{Expected Masking Rate}

\begin{proposition}
Under balanced masking with $\E[i^l] = 0.5$:
\begin{equation}
    \E_l[g_{\text{bal}}(i^l, t)] = 0.5 \quad \forall t \in [0,1]
\end{equation}
\end{proposition}

\begin{proof}
\begin{align}
    \E_l[g_{\text{bal}}(i^l, t)] &= t \cdot \E_l[g_{\text{inv}}(i^l)] + (1{-}t) \cdot \E_l[g_{\text{prop}}(i^l)] \\
    &= t \cdot (\eta + (1{-}2\eta)(1{-}0.5)) + (1{-}t) \cdot (\eta + (1{-}2\eta) \cdot 0.5) \\
    &= t \cdot 0.5 + (1{-}t) \cdot 0.5 = 0.5 \qedhere
\end{align}
\end{proof}

\subsection{Curriculum Interpretation}

The balanced schedule implements a two-phase curriculum that varies smoothly with diffusion timestep:

\begin{table}[h]
\centering
\caption{Curriculum phases in balanced masking.}
\begin{tabular}{lccl}
\toprule
\textbf{Phase} & \textbf{Timestep} & \textbf{Dominant Term} & \textbf{Effect on Important Tokens} \\
\midrule
Early denoising & $t \to 1$ & $g_{\text{inv}}$ & Preserved as semantic anchors \\
Mid denoising & $t = 0.5$ & Equal mix & Balanced masking/preservation \\
Late denoising & $t \to 0$ & $g_{\text{prop}}$ & Focused on for training signal \\
\bottomrule
\end{tabular}
\end{table}

\textbf{Intuition}: When most of the sequence is corrupted ($t \to 1$), the model needs structural guidance---important tokens serve as ``landmarks'' for reconstruction. When most tokens are already recovered ($t \to 0$), the model should practice on the hardest tokens, which are precisely the important ones.

\subsection{Comparison with Static Strategies}

\begin{table}[h]
\centering
\caption{Theoretical comparison of masking strategies.}
\begin{tabular}{lccc}
\toprule
\textbf{Strategy} & \textbf{Anchor Preservation} & \textbf{Hard-Token Training} & \textbf{Adaptive} \\
\midrule
Uniform & None & Equal & No \\
Importance-proportional & None & Strong & No \\
Importance-inverse & Strong & None & No \\
ADLM (binary anchor) & Binary & None & No \\
\textbf{Balanced (ours)} & \textbf{At $t{\to}1$} & \textbf{At $t{\to}0$} & \textbf{Yes} \\
\bottomrule
\end{tabular}
\end{table}

\section{Uncertainty-Guided Sampling}
\label{app:uncertainty}

This section provides a detailed treatment of the uncertainty-guided denoising order used during ATAT inference. While training uses the balanced masking curriculum (deterministic given $i^l$ and $t$), inference uses a \emph{dynamic} unmasking order that adapts to the model's current predictions at each step.

\subsection{Uncertainty Metrics}

We define three uncertainty metrics and analyze their properties. Let $p^l_k = \softmax(\mathbf{o}^l)_k$ denote the denoiser's predicted probability that position $l$ contains token $k$.

\paragraph{Predictive Entropy.}
\begin{equation}
    u^l_{\text{entropy}} = -\sum_{k=1}^K p^l_k \log(p^l_k + \epsilon), \quad \epsilon = 10^{-8}
\end{equation}
Range: $[0, \log K] = [0, 10.83]$ for $K = 50{,}257$. Maximum at the uniform distribution over the vocabulary (maximum ignorance); minimum at a point mass on a single token (maximum confidence). Entropy is our default metric because it is the unique measure satisfying the Shannon axioms: it is continuous, monotonically increasing in the number of likely tokens, and additive for independent predictions. The $\epsilon$ term prevents numerical issues with log(0).

\paragraph{Predictive Variance.}
\begin{equation}
    u^l_{\text{var}} = \frac{1}{K} \sum_{k=1}^K (p^l_k - 1/K)^2
\end{equation}
Range: $[0, (K{-}1)/K^2]$. Measures deviation from the uniform distribution. Unlike entropy, variance is symmetric around the uniform distribution but does not satisfy the Shannon axioms. In practice, variance correlates strongly with entropy ($\rho > 0.95$ in our experiments) but provides slightly less discriminative power for positions with medium uncertainty.

\paragraph{Inverse Max Probability.}
\begin{equation}
    u^l_{\text{maxprob}} = 1 - \max_k p^l_k
\end{equation}
Range: $[0, (K{-}1)/K]$. A computationally efficient alternative to entropy that considers only the mode of the distribution. It is $O(K)$ to compute (single max operation) vs.\ $O(K \log K)$ for entropy (element-wise log and sum). However, it discards information about the shape of the distribution: two positions with $\max_k p_k = 0.5$ but different distributional shapes (one concentrated on 2 tokens, one spread across 10) would receive the same uncertainty score under this metric but different entropy scores.

\subsection{Combined Priority Score}

The denoising priority for each masked position is:
\begin{equation}
    s^l = u^l \cdot i^l
\end{equation}

This product ensures that we unmask positions that are both:
\begin{itemize}[leftmargin=*,itemsep=2pt]
    \item \textbf{Important} ($i^l$ high): The token carries significant information content. Accurately reconstructing this position has a large impact on the overall sequence quality.
    \item \textbf{Uncertain} ($u^l$ high): The model is not yet confident about this position. Unmasking it provides maximal new information to guide subsequent predictions at neighboring positions.
\end{itemize}

Positions with high importance but low uncertainty are already well-predicted and can wait---they will likely be correctly predicted in a single step whenever they are reached. Positions with low importance are not worth prioritizing regardless of uncertainty, as their reconstruction provides minimal contextual benefit to remaining positions.

\paragraph{Alternative Priority Functions.} We considered several alternatives to the multiplicative priority $s^l = u^l \cdot i^l$:
\begin{itemize}[leftmargin=*,itemsep=1pt]
    \item \textbf{Additive:} $s^l = \alpha \cdot u^l + (1 - \alpha) \cdot i^l$ --- Requires an additional hyperparameter $\alpha$ and does not enforce the ``both high'' constraint (high $u$ alone could trigger early unmasking).
    \item \textbf{Minimum:} $s^l = \min(u^l, i^l)$ --- Too conservative; many positions have moderate scores on both axes but low minimum.
    \item \textbf{Importance-only:} $s^l = i^l$ --- Ignores the model's current confidence, leading to a fixed unmasking order independent of the partially reconstructed sequence.
\end{itemize}
In preliminary experiments, the multiplicative formulation consistently performed best (see Table~\ref{tab:denoising-strategies} for a comparison), confirming that the \emph{product} naturally captures the desired ``high on both axes'' semantics.

\paragraph{Unmasking Schedule.} Given the priority scores $\{s^l\}$ for all currently masked positions at diffusion step $t$, we unmask the top-$k$ positions, where $k$ is determined by the noise schedule:
\begin{equation}
    k(t) = \lfloor L \cdot |\alpha_{t+\Delta t} - \alpha_t| \rfloor
\end{equation}
At each step, the number of tokens to unmask is proportional to the change in the noise parameter $\alpha_t$. For a log-linear schedule with uniform time spacing ($\Delta t = 1/\text{NFE}$), this gives approximately $k = L / \text{NFE}$ tokens per step ($\approx 1$ token per step at NFE$=$1000 with $L = 1024$).

\section{Extended Experimental Results}
\label{app:experiments}

This section provides comprehensive experimental results beyond those presented in the main paper, including full ablation sweeps, per-domain breakdowns, training dynamics, and additional analyses.

\subsection{Full Zero-Shot Perplexity Breakdown}

\begin{table}[h]
\centering
\caption{Complete zero-shot perplexity results across all 7 benchmarks with standard deviations over 3 evaluation runs. All models trained only on OpenWebText. Best DLM result in \textbf{bold}, best overall in \underline{underline}.}
\label{tab:zero-shot-full}
\begin{tabular}{lccccccc|c}
\toprule
\textbf{Model} & \textbf{LAM.} & \textbf{PTB} & \textbf{Wiki.} & \textbf{LM1B} & \textbf{AG} & \textbf{PubMed} & \textbf{ArXiv} & \textbf{Avg.} \\
\midrule
AR (GPT-2) & 51.28 & \underline{82.05} & \underline{25.75} & \underline{21.55} & 52.09 & 49.01 & 41.73 & 46.21 \\
\midrule
D3PM & 65.41 & 121.30 & 47.50 & 47.50 & 78.92 & 62.17 & 55.83 & 68.38 \\
SEDD & 49.86 & 100.09 & 34.28 & 32.79 & 65.34 & 45.22 & 40.51 & 52.58 \\
MDLM & 47.52 & 95.26 & 32.83 & 27.04 & 61.15 & 41.89 & 37.37 & 49.01 \\
ADLM & 44.32 & 95.37 & 31.94 & 24.46 & 55.72 & 37.56 & 33.69 & 46.15 \\
\midrule
ATAT (500K) & 44.97{\scriptsize$\pm$0.21} & 95.48{\scriptsize$\pm$0.34} & 32.15{\scriptsize$\pm$0.18} & 25.89{\scriptsize$\pm$0.15} & 55.03{\scriptsize$\pm$0.28} & 37.21{\scriptsize$\pm$0.23} & 33.42{\scriptsize$\pm$0.19} & 46.31 \\
ATAT (1M) & \textbf{43.52}{\scriptsize$\pm$0.18} & \textbf{94.63}{\scriptsize$\pm$0.29} & \textbf{30.47}{\scriptsize$\pm$0.14} & \textbf{24.21}{\scriptsize$\pm$0.12} & \textbf{53.49}{\scriptsize$\pm$0.24} & \textbf{35.97}{\scriptsize$\pm$0.20} & \textbf{32.08}{\scriptsize$\pm$0.16} & \textbf{44.91} \\
\bottomrule
\end{tabular}
\end{table}

\subsection{Per-Domain Improvement Analysis}

Table~\ref{tab:domain-improvement} analyzes ATAT's improvements relative to both the AR baseline and the best prior DLM (ADLM) on each benchmark, decomposing gains into absolute and relative terms.

\begin{table}[h]
\centering
\caption{Per-domain improvement analysis: ATAT (1M) vs.\ AR baseline and ADLM. Negative values indicate ATAT is better (lower perplexity).}
\label{tab:domain-improvement}
\begin{tabular}{lcccccc}
\toprule
\textbf{Domain} & \textbf{AR} & \textbf{ADLM} & \textbf{ATAT} & \textbf{$\Delta$ vs AR} & \textbf{$\Delta$ vs ADLM} & \textbf{\% vs AR} \\
\midrule
LAMBADA & 51.28 & 44.32 & 43.52 & $-$7.76 & $-$0.80 & $-$15.1\% \\
PTB & 82.05 & 95.37 & 94.63 & +12.58 & $-$0.74 & +15.3\% \\
WikiText-2 & 25.75 & 31.94 & 30.47 & +4.72 & $-$1.47 & +18.3\% \\
LM1B & 21.55 & 24.46 & 24.21 & +2.66 & $-$0.25 & +12.3\% \\
AG News & 52.09 & 55.72 & 53.49 & +1.40 & $-$2.23 & +2.7\% \\
PubMed & 49.01 & 37.56 & 35.97 & $-$13.04 & $-$1.59 & $-$26.6\% \\
ArXiv & 41.73 & 33.69 & 32.08 & $-$9.65 & $-$1.61 & $-$23.1\% \\
\midrule
\textbf{Average} & 46.21 & 46.15 & 44.91 & $-$1.30 & $-$1.24 & $-$2.8\% \\
\bottomrule
\end{tabular}
\end{table}

\paragraph{Key Observations.} (1)~ATAT beats AR on 3 of 7 benchmarks (LAMBADA, PubMed, ArXiv), with losses on PTB, WikiText-2, LM1B, and AG News where the left-to-right factorization closely matches the sequential nature of the evaluation corpora. (2)~ATAT's largest advantages over AR come on specialized domains (PubMed: $-$26.6\%, ArXiv: $-$23.1\%) where learned importance identifies domain-specific terminology. (3)~Relative to ADLM, ATAT improves on every benchmark, demonstrating that learned hybrid importance strictly dominates frequency-based anchoring. The largest gains are on AG News ($-$2.23), PubMed ($-$1.59), and ArXiv ($-$1.61), where contextual importance disambiguates tokens that frequency treats uniformly.

\subsection{Training Checkpoint Progression}

\begin{table}[h]
\centering
\caption{ATAT performance progression across training checkpoints (WikiText-2 validation perplexity and importance estimator oracle correlation).}
\label{tab:checkpoint-progression}
\begin{tabular}{rccccl}
\toprule
\textbf{Steps} & \textbf{Tokens} & \textbf{Wiki PPL} & \textbf{LM1B PPL} & \textbf{Oracle Corr.} & \textbf{Notes} \\
\midrule
10K & 0.66B & 58.42 & 42.17 & 0.58 & Early training, estimator learning \\
50K & 3.28B & 42.31 & 34.89 & 0.75 & Ablation checkpoint \\
100K & 6.55B & 38.91 & 31.52 & 0.79 & Estimator near convergence \\
200K & 13.1B & 35.67 & 28.94 & 0.80 & Denoiser primary learner \\
500K & 32.8B & 32.15 & 25.89 & 0.81 & Reported 500K checkpoint \\
750K & 49.2B & 31.18 & 24.89 & 0.82 & Diminishing returns \\
1M & 65.5B & 30.47 & 24.21 & 0.82 & Final reported checkpoint \\
\bottomrule
\end{tabular}
\end{table}

The checkpoint progression reveals two distinct learning regimes: (1)~\textit{Rapid importance learning} (0--100K steps): the oracle correlation rises from 0.0 to 0.79 as the importance MLP learns to extract surprisal signals from GPT-2 representations. During this phase, perplexity drops steeply as the model transitions from uniform to informed masking. (2)~\textit{Denoiser refinement} (100K--1M steps): the oracle correlation saturates near 0.82 while perplexity continues to decrease, indicating that the denoiser is learning to better exploit the now-stable importance signal.

\subsection{Full Ablation: Importance Estimation}

\begin{table}[h]
\centering
\caption{Importance estimation ablation on WikiText-2 validation (50K-step runs, 10\% OpenWebText subset).}
\label{tab:imp-ablation-full}
\begin{tabular}{lccc}
\toprule
\textbf{Strategy} & \textbf{PPL} $\downarrow$ & \textbf{$\Delta$ vs Uniform} & \textbf{Oracle Corr.} \\
\midrule
Uniform (no importance) & 42.31 & --- & --- \\
Frequency-only ($\lambda=0$) & 41.87 & $-$0.44 & 0.71 \\
Context-only ($\lambda=1$) & 40.12 & $-$2.19 & 0.78 \\
Hybrid ($\lambda=0.5$) & 39.56 & $-$2.75 & 0.79 \\
Hybrid ($\lambda=0.6$) & 39.21 & $-$3.10 & 0.80 \\
\textbf{Hybrid ($\lambda=0.7$)} & \textbf{39.03} & $\mathbf{-3.28}$ & \textbf{0.82} \\
Hybrid ($\lambda=0.8$) & 39.15 & $-$3.16 & 0.81 \\
Hybrid ($\lambda=0.9$) & 39.87 & $-$2.44 & 0.80 \\
\bottomrule
\end{tabular}
\end{table}

\paragraph{Detailed Analysis.} Several patterns emerge from the importance estimation ablation:

\begin{enumerate}[leftmargin=*,itemsep=2pt]
    \item \textbf{Frequency alone provides limited benefit} ($-$0.44 PPL). While inverse frequency is a useful prior (rare tokens are typically more important), it is context-independent: the word ``cell'' receives the same frequency-based importance whether it appears in a biology paper (high-importance) or as ``prison cell'' (moderate-importance). The oracle correlation of 0.71 confirms that frequency captures only a coarse approximation of true importance.
    
    \item \textbf{Context alone is substantially better} ($-$2.19 PPL). GPT-2 surprisal captures token-level predictability conditioned on the full left context, yielding oracle correlation 0.78. However, context-only importance is noisy for rare tokens where GPT-2 has limited data---the estimator may assign inconsistent importance to infrequent technical terms.
    
    \item \textbf{The hybrid combination is synergistic, not merely additive.} If the improvements were additive, we would expect $-$0.44 $+$ ($-$2.19) $=$ $-$2.63 PPL improvement. The best hybrid ($\lambda = 0.7$) achieves $-$3.28 PPL, an additional 0.65 PPL beyond additivity. This \emph{synergy} arises because frequency and context importance are \emph{complementary}: frequency stabilizes the importance estimate for rare tokens (where context is noisy), while context provides fine-grained discrimination for common tokens (where frequency provides little signal).
    
    \item \textbf{The optimal $\lambda = 0.7$ assigns more weight to context}, reflecting that contextual information is more discriminative. However, the 30\% frequency component is not negligible: removing it ($\lambda = 1.0$ vs.\ $\lambda = 0.7$) costs 1.09 PPL, confirming the stabilizing role of frequency for rare tokens.
    
    \item \textbf{The oracle correlation increases monotonically up to $\lambda = 0.7$} and then plateaus. This suggests that $\lambda = 0.7$ achieves near-optimal alignment with the oracle importance (GPT-2 full-sequence surprisal), and further increases in $\lambda$ introduce noise without improving alignment.
\end{enumerate}

\subsection{Full Ablation: Masking Strategies}

\begin{table}[h]
\centering
\caption{Masking strategy ablation (50K-step runs). PPL on WikiText-2 validation.}
\label{tab:mask-ablation-full}
\begin{tabular}{lcc}
\toprule
\textbf{Strategy} & \textbf{PPL} $\downarrow$ & \textbf{Notes} \\
\midrule
Uniform & 42.31 & Standard DLM baseline \\
Importance-proportional only & 41.12 & More training signal on hard tokens \\
Importance-inverse only & 40.89 & Anchor preservation helps more \\
\textbf{Balanced} & \textbf{39.03} & Combines both via time interpolation \\
\bottomrule
\end{tabular}
\end{table}

\paragraph{Detailed Analysis.} The masking strategy ablation reveals an important design principle for importance-aware diffusion:

\begin{itemize}[leftmargin=*,itemsep=2pt]
    \item \textbf{Importance-proportional masking} masks high-importance tokens more frequently, forcing the denoiser to practice on the hardest reconstructions. This provides a 1.19 PPL improvement over uniform but is limited because it also removes the most informative context tokens, reducing the denoiser's ability to leverage context.
    
    \item \textbf{Importance-inverse masking} does the opposite: it preferentially preserves high-importance tokens as context, allowing the denoiser to reconstruct easy (low-importance) tokens with strong contextual support. Surprisingly, this outperforms proportional masking by 0.23 PPL, suggesting that \emph{context preservation is more valuable than hard-example training} in the masked diffusion setting.
    
    \item \textbf{Balanced masking} interpolates between these strategies over diffusion time: at early $t$ (few masked tokens), proportional masking dominates, ensuring the denoiser gets hard examples; at late $t$ (many masked tokens), inverse masking dominates, preserving context. This time-dependent interpolation captures the best of both strategies and achieves a 3.28 PPL improvement, substantially exceeding either component alone.
    
    \item \textbf{The gap between inverse-only and balanced} (1.86 PPL) demonstrates that the time-dependent interpolation is not merely cosmetic---it captures a fundamental property of the diffusion process: the optimal masking strategy changes as the noise level changes.
\end{itemize}

\subsection{NFE Trade-off}

\begin{table}[h]
\centering
\caption{Number of Function Evaluations (NFE) vs.\ perplexity at generation time. Both models evaluated on WikiText-2 test (1M-step checkpoints).}
\label{tab:nfe-full}
\begin{tabular}{lcccc}
\toprule
\textbf{NFE} & \textbf{ATAT PPL} & \textbf{MDLM PPL} & \textbf{$\Delta$ (abs.)} & \textbf{$\Delta$\%} \\
\midrule
50 & 51.8 & 56.2 & $-$4.4 & $-$7.8\% \\
100 & 45.2 & 48.1 & $-$2.9 & $-$6.0\% \\
250 & 41.8 & 44.5 & $-$2.7 & $-$6.1\% \\
500 & 40.1 & 42.8 & $-$2.7 & $-$6.3\% \\
1000 & 39.0 & 41.2 & $-$2.2 & $-$5.3\% \\
\bottomrule
\end{tabular}
\end{table}

\paragraph{Detailed Analysis.} ATAT consistently outperforms MDLM across all NFE budgets, with the largest \emph{absolute} advantage at the lowest NFE (4.4 PPL at 50 NFE, 2.9 at 100 NFE). This pattern has a clear explanation:

\begin{itemize}[leftmargin=*,itemsep=2pt]
    \item At low NFE, each denoising step must be maximally efficient. ATAT's uncertainty-guided denoising order ensures that the few available steps are spent on the most impactful positions---reconstructing structural tokens first to build context, then using that context for content token prediction. Random order (MDLM) wastes steps on positions that could have been predicted from context.
    
    \item At high NFE (1000), both models have enough budget to reconstruct most tokens accurately regardless of order. The remaining gap (2.2 PPL) reflects the training-time advantage of balanced masking: ATAT's denoiser has learned to be a better reconstructor overall, independent of the inference strategy.
    
    \item The 50 NFE setting (4.4 PPL gap) is particularly relevant for practical applications where inference latency is constrained. ATAT at 50 NFE (51.8 PPL) outperforms MDLM at 100 NFE (48.1 PPL... note ATAT at 100 NFE is 45.2 PPL), achieving $2\times$ better inference efficiency for comparable quality.
\end{itemize}

\subsection{Smoothing Parameter $\eta$ Sensitivity}

\begin{table}[h]
\centering
\caption{Sensitivity to smoothing parameter $\eta$ (50K-step runs, WikiText-2 validation). The smoothing parameter interpolates masking probabilities toward 0.5.}
\label{tab:eta-sensitivity}
\begin{tabular}{ccl}
\toprule
$\eta$ & \textbf{PPL} $\downarrow$ & \textbf{Interpretation} \\
\midrule
0.0 & 40.15 & No smoothing; some tokens deterministically masked/preserved \\
0.1 & 39.42 & Mild smoothing; still strong importance differentiation \\
0.2 & 39.18 & Moderate smoothing \\
\textbf{0.3} & \textbf{39.03} & Optimal exploration-exploitation balance \\
0.4 & 39.31 & Over-smoothed; importance signal attenuated \\
0.5 & 42.31 & Fully uniform; $g(i, t) = 0.5$ for all tokens \\
\bottomrule
\end{tabular}
\end{table}

\paragraph{Detailed Analysis.} The smoothing parameter $\eta$ controls the trade-off between exploitation (faithful use of learned importance for masking) and exploration (ensuring all tokens have nonzero probability of being masked or preserved):

\begin{itemize}[leftmargin=*,itemsep=2pt]
    \item \textbf{$\eta = 0$} (no smoothing): The masking probability is a deterministic function of importance, meaning some high-importance tokens are always masked at certain timesteps. This creates ``dead zones'' where the denoiser never sees certain tokens in context, limiting generalization. The 1.12 PPL degradation vs.\ optimal confirms this pathology.
    
    \item \textbf{$\eta = 0.3$} (optimal): Every token has at least a 15\% probability of being in any masking state (the minimum masking probability is $\eta/2 = 0.15$). This ensures the denoiser sees diverse masking patterns while still strongly differentiating between high- and low-importance tokens.
    
    \item \textbf{$\eta = 0.5$}: The masking probability becomes $g(i, t) = 0.5$ for all tokens regardless of importance, recovering the uniform baseline. The fact that PPL reverts exactly to 42.31 (the uniform baseline) confirms that our smoothing implementation is correct and that $\eta = 0.5$ fully neutralizes importance information.
    
    \item The transition from $\eta = 0.3$ to $\eta = 0.5$ (3.28 PPL) is much steeper than from $\eta = 0.0$ to $\eta = 0.3$ (1.12 PPL), suggesting that the importance signal provides substantially more value than the exploration benefit of smoothing.
\end{itemize}

\subsection{Importance Weight $\beta$ and Supervision Weight $\gamma$}

\begin{table}[h]
\centering
\caption{Sensitivity to loss hyperparameters (50K-step runs, WikiText-2 validation). $\beta$ controls importance-weighting in the denoising loss; $\gamma$ controls the importance supervision loss magnitude.}
\label{tab:beta-gamma}
\begin{tabular}{cccc}
\toprule
$\beta$ & $\gamma$ & \textbf{PPL} $\downarrow$ & \textbf{Interpretation} \\
\midrule
0.0 & 0.003 & 40.45 & No importance weighting in loss (uniform loss weights) \\
0.5 & 0.003 & 39.52 & Mild importance weighting \\
\textbf{1.0} & \textbf{0.003} & \textbf{39.03} & Selected: equal weighting of importance \\
2.0 & 0.003 & 39.28 & Over-weighting; high-importance tokens dominate \\
\midrule
1.0 & 0.0 & 40.89 & No importance supervision (unsupervised estimator) \\
1.0 & 0.001 & 39.45 & Weak supervision; slow estimator convergence \\
\textbf{1.0} & \textbf{0.003} & \textbf{39.03} & Selected \\
1.0 & 0.01 & 39.67 & Over-regularized; supervision dominates training \\
\bottomrule
\end{tabular}
\end{table}

\paragraph{Detailed Analysis.}

\begin{itemize}[leftmargin=*,itemsep=2pt]
    \item \textbf{$\beta$ analysis:} The importance weight $\beta$ scales the per-token loss contribution by $1 + \beta \cdot i^l$. At $\beta = 0$, all tokens contribute equally; at $\beta = 1$, a token with $i^l = 0.8$ contributes $1.8\times$ the loss of a token with $i^l = 0$. The optimal $\beta = 1.0$ doubles the gradient signal from the most important tokens (those with $i^l \approx 1$), providing stronger supervision on exactly the tokens that matter most for perplexity. At $\beta = 2.0$, the model over-focuses on high-importance tokens and under-trains on function words, leading to grammatical disfluencies despite better content word prediction.
    
    \item \textbf{$\gamma$ analysis:} The supervision coefficient $\gamma$ must be carefully balanced. At $\gamma = 0$ (no supervision), the importance estimator receives no direct signal and can only learn importance indirectly through the denoising loss gradient, resulting in slow convergence and poor importance estimates (PPL 40.89). At $\gamma = 0.01$, the supervision loss becomes ${\sim}$10\% of the total loss, causing the importance estimator to ``overfit'' to oracle targets at the expense of the denoiser's needs. The optimal $\gamma = 0.003$ ensures supervision contributes ${\sim}$1--3\% of the total gradient magnitude, providing a useful learning signal without distorting the primary denoising objective.
\end{itemize}

\subsection{POS-Tag Importance Distribution}

\begin{table}[h]
\centering
\caption{Mean learned importance by part-of-speech category (WikiText-2 validation set, ATAT 1M checkpoint). POS tags assigned using spaCy's \texttt{en\_core\_web\_sm} model.}
\label{tab:pos-full}
\begin{tabular}{lcccc}
\toprule
\textbf{POS Category} & \textbf{Mean $i^l$} & \textbf{Std Dev} & \textbf{Median $i^l$} & \textbf{Expected Role} \\
\midrule
Proper nouns (NNP) & 0.78 & 0.12 & 0.81 & Content (entities) \\
Common nouns (NN) & 0.62 & 0.15 & 0.64 & Content (concepts) \\
Verbs (VB*) & 0.58 & 0.14 & 0.59 & Content (actions) \\
Adjectives (JJ) & 0.54 & 0.16 & 0.55 & Content (modifiers) \\
Adverbs (RB) & 0.48 & 0.18 & 0.47 & Mixed \\
Prepositions (IN) & 0.32 & 0.11 & 0.31 & Function \\
Determiners (DT) & 0.21 & 0.09 & 0.20 & Function \\
Punctuation & 0.15 & 0.08 & 0.13 & Structural \\
\bottomrule
\end{tabular}
\end{table}

\paragraph{Detailed Analysis.} The learned importance distribution reveals linguistically meaningful patterns:

\begin{enumerate}[leftmargin=*,itemsep=2pt]
    \item \textbf{Clear content-function dichotomy:} Content words (NNP, NN, VB, JJ) cluster in the $[0.54, 0.78]$ range, while function words (IN, DT) and punctuation cluster in $[0.15, 0.32]$. The gap between the lowest content category (adjectives, 0.54) and the highest function category (prepositions, 0.32) is 0.22, demonstrating strong discrimination.
    
    \item \textbf{Within-category variance is informative:} The standard deviations (0.08--0.18) are substantial relative to the means, indicating that POS alone does not determine importance. For example, the adjective ``novel'' in a scientific context receives $i^l \approx 0.72$ (well above the JJ mean of 0.54), while the adjective ``big'' receives $i^l \approx 0.38$ because it is highly predictable from context. This within-POS variation is captured by the context component of hybrid importance.
    
    \item \textbf{Adverbs are the most heterogeneous category} (std dev 0.18), reflecting their dual role: some adverbs are content-bearing (``significantly'', ``dramatically'': $i^l \approx 0.65$) while others are functional (``very'', ``also'': $i^l \approx 0.28$).
    
    \item \textbf{Proper nouns have the highest mean importance} (0.78), consistent with their role as entity names that are low-frequency and context-dependent. The std dev of 0.12 is moderate: well-known proper nouns (``the United States'': $i^l \approx 0.62$) receive lower importance than rare ones (``Schwarzenegger'': $i^l \approx 0.91$).
    
    \item \textbf{Comparison with frequency-only importance:} Under a pure frequency-based scheme, determiners would all receive similarly low importance. Under ATAT's hybrid scheme, the determiner ``the'' receives $i^l = 0.08$ (extremely common), while ``those'' receives $i^l = 0.34$ (less common, more context-dependent). This demonstrates the value of the learned contextual component.
\end{enumerate}

\subsection{Training Dynamics and Convergence Analysis}

The importance estimator and denoiser exhibit markedly different convergence behaviors, which we analyze in detail.

\begin{table}[h]
\centering
\caption{Detailed training dynamics: denoising loss ($\gL_{\text{denoise}}$), importance supervision loss ($\gL_{\text{importance}}$), and oracle correlation at key training checkpoints. The importance estimator converges ${\sim}5\times$ faster than the denoiser.}
\label{tab:training-dynamics}
\begin{tabular}{rcccccc}
\toprule
\textbf{Step} & $\gL_{\text{denoise}}$ & $\gL_{\text{importance}}$ & \textbf{Oracle Corr.} & \textbf{Wiki PPL} & \textbf{GPU Mem.} & \textbf{Throughput (per GPU)} \\
\midrule
0 & 10.82 & 0.248 & 0.00 & --- & 19.8 GB & --- \\
1K & 8.45 & 0.142 & 0.31 & --- & 19.8 GB & 63K tok/s \\
5K & 6.92 & 0.071 & 0.52 & 72.41 & 19.8 GB & 63K tok/s \\
10K & 6.21 & 0.045 & 0.58 & 58.42 & 19.8 GB & 63K tok/s \\
20K & 5.73 & 0.028 & 0.75 & 48.17 & 19.8 GB & 63K tok/s \\
50K & 5.18 & 0.015 & 0.75 & 42.31 & 19.8 GB & 63K tok/s \\
100K & 4.82 & 0.009 & 0.79 & 38.91 & 19.8 GB & 63K tok/s \\
200K & 4.51 & 0.006 & 0.80 & 35.67 & 19.8 GB & 63K tok/s \\
500K & 4.15 & 0.004 & 0.81 & 32.15 & 19.8 GB & 63K tok/s \\
1M & 3.89 & 0.003 & 0.82 & 30.47 & 19.8 GB & 63K tok/s \\
\bottomrule
\end{tabular}
\end{table}

\paragraph{Importance Estimator Convergence.} The importance MLP converges rapidly:
\begin{itemize}[leftmargin=*,itemsep=1pt]
    \item \textbf{0--20K steps}: Oracle correlation rises from 0.0 to 0.75, accounting for 91\% of the total correlation improvement. The MSE loss drops from 0.248 to 0.028 ($8.9\times$ reduction). This rapid convergence is expected given the small model size (200K parameters) and direct supervision from GPT-2 surprisal targets.
    \item \textbf{20K--100K steps}: Correlation increases from 0.75 to 0.79, with diminishing returns. The MLP refines its predictions on ambiguous tokens where surprisal is context-dependent.
    \item \textbf{100K--1M steps}: Near-saturation, with correlation increasing from 0.79 to 0.82. The remaining improvements come from better handling of rare tokens and boundary cases.
\end{itemize}

\paragraph{Denoiser Convergence.} The denoiser (48M parameters) follows a different trajectory:
\begin{itemize}[leftmargin=*,itemsep=1pt]
    \item The denoising loss decreases steadily from 10.82 to 3.89 across 1M steps, with no plateau observed, suggesting that further training could yield additional improvements.
    \item The denoiser benefits from the early convergence of the importance estimator: once masking becomes importance-aware (by ${\sim}$20K steps), the denoiser receives a more informative training signal at every subsequent step.
    \item The decoupling of convergence rates validates our architectural choice to separate importance estimation from denoising: the lightweight MLP quickly converges to provide stable importance signals, allowing the larger denoiser to focus entirely on the reconstruction task.
\end{itemize}

\paragraph{Memory and Throughput.} Training requires 19.8 GB per GPU (peak, with mixed-precision FP16) across all 4$\times$ RTX 4090 24GB GPUs. At batch size 64 and sequence length 1024, each training step processes 65,536 tokens, requiring ${\sim}$0.26 seconds per step. The total throughput is ${\sim}$253K tokens/second across the 4-GPU setup (${\sim}$63K tokens/second per GPU). The full 1M-step training run thus requires ${\sim}$72 hours of wall-clock time.

\section{Implementation Details}
\label{app:implementation}

This section provides comprehensive implementation details to support reproducibility of our results.

\subsection{Architecture Specifications}

\begin{table}[H]
\centering
\caption{ATAT architecture specifications.}
\label{tab:arch-specs}
\begin{tabular}{lcc}
\toprule
\textbf{Component} & \textbf{Configuration} & \textbf{Parameters} \\
\midrule
\multicolumn{3}{l}{\textit{Anchor Network (Frozen)}} \\
Model & GPT-2 Small & 124M \\
Layers & 12 & --- \\
Hidden dimension & 768 & --- \\
Attention heads & 12 & --- \\
Vocabulary size & 50,257 & --- \\
Max sequence length & 1024 & --- \\
Position embeddings & Learned absolute & --- \\
Weight source & HuggingFace \texttt{gpt2} & --- \\
\midrule
\multicolumn{3}{l}{\textit{Importance Estimator (Trainable)}} \\
Architecture & 2-layer MLP & 200K \\
Input dimension & 768 & --- \\
Hidden dimension & 256 & --- \\
Output dimension & 1 & --- \\
Activation & GELU & --- \\
Normalization & LayerNorm (pre-activation) & --- \\
Output activation & Sigmoid ($\to [0,1]$) & --- \\
\midrule
\multicolumn{3}{l}{\textit{Denoiser (Trainable)}} \\
Architecture & 6-layer Transformer & 48M \\
Hidden dimension & 768 & --- \\
Attention heads & 12 & --- \\
Head dimension & 64 & --- \\
Feed-forward dimension & 3072 & --- \\
Feed-forward activation & GELU & --- \\
Dropout & 0.1 & --- \\
Attention type & Full bidirectional & --- \\
Position encoding & Rotary (RoPE) & --- \\
Normalization & Pre-LayerNorm & --- \\
\midrule
\multicolumn{3}{l}{\textit{Importance Projection}} \\
Projection $\mathbf{W}_i$ & $\R^{768 \times 2}$ & 1,536 \\
Input format & $[i^l;\; 1{-}i^l]$ & --- \\
\midrule
\textbf{Total Trainable} & & \textbf{$\sim$49M} \\
\textbf{Total (incl.\ frozen)} & & \textbf{$\sim$173M} \\
\bottomrule
\end{tabular}
\end{table}

\paragraph{Detailed Parameter Accounting.} The importance estimator consists of a LayerNorm ($768 \times 2 = 1{,}536$ params), a first linear layer ($768 \times 256 + 256 = 196{,}864$), and a second linear layer ($256 \times 1 + 1 = 257$), totaling $198{,}657$ trainable parameters. Including the importance projection $\mathbf{W}_i \in \R^{768 \times 2}$ ($2 \times 768 + 768 = 2{,}304$ params), the full importance pathway totals $\sim$200K parameters. The denoiser's 48M parameters include 6 Transformer blocks, each containing multi-head self-attention ($4 \times 768 \times 768 = 2.36$M), feed-forward network ($768 \times 3072 + 3072 \times 768 = 4.72$M), two LayerNorm layers ($768 \times 4 = 3{,}072$), plus input/output embedding projections.

\subsection{Denoiser Initialization Strategy}

% The denoiser is initialized from the \textbf{first 6 layers} of the frozen GPT-2 Small checkpoint. Specifically:
% \begin{itemize}[itemsep=2pt, topsep=4pt, parsep=0pt]
%     \item Self-attention weights ($\mathbf{W}_Q, \mathbf{W}_K, \mathbf{W}_V, \mathbf{W}_O$) are copied from GPT-2 layers 0--5.
%     \item Feed-forward weights ($\mathbf{W}_1, \mathbf{W}_2$) are copied similarly.
%     \item LayerNorm parameters are copied from the corresponding GPT-2 layers.
%     \item The output projection head $\mathbf{W}_{\text{out}} \in \R^{K \times 768}$ is initialized from GPT-2's language model head (the token embedding matrix, transposed).
%     \item The importance projection $\mathbf{W}_i \in \R^{768 \times 2}$ is initialized with Xavier uniform initialization.
% \end{itemize}

This initialization strategy serves dual purposes: (1)~it provides a warm start that accelerates convergence by approximately $2\times$ compared to random initialization (based on preliminary experiments showing 50K-step random-init achieves only 45.8 PPL vs.\ 42.3 with warm start), and (2)~it ensures the denoiser's representation space is compatible with the frozen anchor's output, since both share the same architectural family. After initialization, all denoiser weights are unfrozen and updated during training; the denoiser specializes for the denoising task and diverges from the original GPT-2 weights.

\subsection{Training Configuration}

\begin{table}[H]
\centering
\caption{Complete training hyperparameters for the final ATAT models.}
\label{tab:training-config}
\begin{tabular}{ll}
\toprule
\textbf{Hyperparameter} & \textbf{Value} \\
\midrule
\multicolumn{2}{l}{\textit{Optimizer}} \\
Algorithm & AdamW~\citep{adamw} \\
Learning rate & $3 \times 10^{-4}$ \\
$\beta_1$, $\beta_2$ & 0.9, 0.999 \\
$\epsilon$ & $1 \times 10^{-8}$ \\
Weight decay & 0.01 \\
LR schedule & Cosine decay to 0 \\
Warmup steps & 1,000 (linear) \\
Gradient clipping & Max norm 1.0 \\
\midrule
\multicolumn{2}{l}{\textit{Data}} \\
Batch size (total) & 64 \\
Batch size (per GPU) & 16 \\
Sequence length & 1024 \\
Total steps & 1,000,000 \\
Effective tokens per step & 65,536 \\
Precision & Mixed (FP16 with FP32 master weights) \\
\midrule
\multicolumn{2}{l}{\textit{ATAT-specific hyperparameters}} \\
$\lambda$ (hybrid weight) & 0.7 \\
$\beta$ (importance loss weight) & 1.0 \\
$\gamma$ (supervision weight) & 0.003 \\
$\eta$ (smoothing parameter) & 0.3 \\
$\tau$ (surprisal temperature) & 10.0 \\
\midrule
\multicolumn{2}{l}{\textit{Curriculum scheduler}} \\
Easy phase ($i \in [0, 0.3]$) & Steps 0--200K (first 20\%) \\
Medium phase ($i \in [0.3, 0.7]$) & Steps 200K--600K (next 40\%) \\
Hard phase ($i \in [0, 1.0]$) & Steps 600K--1M (final 40\%) \\
Phase transition & Smooth cosine interpolation \\
Loss plateau detection & Extend phase if loss decreases $<$1\% per 10K steps \\
\midrule
\multicolumn{2}{l}{\textit{Compute}} \\
Hardware & 4$\times$ NVIDIA RTX 4090 (24GB) \\
Distributed strategy & Data-parallel (PyTorch DDP) \\
Training time (1M steps) & $\sim$72 hours \\
Tokens processed (500K) & 32.8B \\
Tokens processed (1M) & 65.5B \\
Peak GPU memory & 19.8 GB per GPU \\
Throughput & $\sim$253K tokens/sec (total), $\sim$63K per GPU \\
\bottomrule
\end{tabular}
\end{table}

\subsection{Gradient Flow and Optimization Details}

\paragraph{Loss Balancing.} The total loss $\gL_{\text{ATAT}} = \gL_{\text{denoise}} + \gamma \gL_{\text{importance}}$ balances two objectives with very different scales. The denoising loss $\gL_{\text{denoise}}$ is a sum of cross-entropy terms over masked positions, typically in the range $[3.5, 11.0]$. The importance supervision loss $\gL_{\text{importance}}$ is an MSE averaged over all positions, typically in the range $[0.001, 0.25]$. The coefficient $\gamma = 0.003$ was selected to ensure that the importance gradient contributes approximately 1--5\% of the total gradient magnitude, preventing the importance objective from dominating early in training while still providing meaningful supervision. We verified that importance gradients flow only to the estimator $f_\varphi$ (since oracle targets are computed from the frozen GPT-2 and detached from the computation graph), while denoising gradients update both the estimator and denoiser jointly.

\paragraph{Gradient Accumulation.} With a per-GPU batch size of 16 and 4 GPUs, the effective batch size is 64 sequences $\times$ 1024 tokens $=$ 65,536 tokens per step. No gradient accumulation beyond the data-parallel synchronization is used, as this batch size was sufficient for stable training.

\paragraph{Mixed Precision.} We use PyTorch's automatic mixed precision (AMP) with FP16 for forward/backward passes and FP32 for optimizer states and master weights. The loss scaling starts at $2^{16}$ and is dynamically adjusted. We observed no numerical instability issues throughout training, which we attribute to the relatively moderate model size and the stabilizing effect of the frozen anchor.

\subsection{Noise Schedule}

We use the log-linear noise schedule from MDLM~\citep{mdlm}:
\begin{equation}
    \alpha_t = 1 - t, \quad t \in [0, 1]
\end{equation}
which gives a uniform rate of masking over time. The continuous-time NELBO weighting becomes $-\dot\alpha_t/(1{-}\alpha_t) = 1/t$, giving higher weight to later timesteps where reconstruction is more informative. In discrete approximation during training, we sample $t \sim \text{Uniform}[0, 1]$ at each step and compute $\alpha_t = 1 - t$. The expected number of masked tokens at time $t$ is $L \cdot (1 - \alpha_t) = L \cdot t$; at $t = 0.5$, approximately half the tokens are masked.

\paragraph{Comparison with Alternative Schedules.} We briefly compared the log-linear schedule against a cosine schedule ($\alpha_t = \cos(\pi t / 2)$) in preliminary experiments. The log-linear schedule performed marginally better (0.3 PPL on WikiText-2 at 50K steps), likely because its uniform masking rate provides a more even distribution of training signal across diffusion timesteps.

\subsection{Data Processing Pipeline}

\paragraph{Dataset.} OpenWebText~\citep{owt} is a public reproduction of the WebText corpus used to train GPT-2. It contains approximately 8 million web documents (38 GB of text) scraped from URLs shared on Reddit with at least 3 upvotes. The corpus covers diverse topics including news, discussions, creative writing, and technical content.

\paragraph{Tokenization.} All text is tokenized using GPT-2's byte-pair encoding (BPE) tokenizer with a vocabulary of $K = 50{,}257$ tokens. The tokenizer handles all Unicode characters through byte-level encoding, ensuring complete coverage of the input text without unknown tokens.

\paragraph{Sequence Packing.} Following MDLM~\citep{mdlm}, we concatenate all documents with $\langle \text{EOS} \rangle$ token separators (token ID 50256) into a single stream, then split into contiguous chunks of length $L = 1024$. This packing strategy:
\begin{itemize}[leftmargin=*,itemsep=1pt]
    \item Avoids padding waste (no tokens are discarded or padded).
    \item Ensures every training sequence has exactly 1024 tokens.
    \item May result in sequences that span document boundaries; the $\langle \text{EOS} \rangle$ token serves as a separator.
    \item Produces approximately 37M training sequences from the full corpus.
\end{itemize}

\paragraph{Frequency Computation.} Token frequencies are computed in a single pass over the tokenized training corpus before training begins. The resulting frequency table maps each token ID $v \in \{0, \ldots, 50{,}256\}$ to its corpus count $\text{freq}(v)$. The most frequent token (``\textvisiblespace the'', token ID 262) has frequency ${\sim}$54M; approximately 3,200 tokens are hapax legomena (frequency 1). The frequency table is stored as a 50,257-dimensional tensor and loaded at the start of each training run. Computation of $i^l_{\text{freq}}$ for an entire batch requires only a single lookup and element-wise operations, adding negligible overhead ($<$0.1\% of step time).

\subsection{Evaluation Protocol}

\paragraph{Perplexity Computation.} Perplexity is computed using the NELBO bound following MDLM~\citep{mdlm}:
\begin{equation}
    \text{PPL} = \exp\left(\frac{\gL_{\text{NELBO}}(\rvx;\theta)}{L}\right)
\end{equation}
where the NELBO is averaged over all $L$ tokens in each test sequence. The NELBO integral over $t \in [0,1]$ is approximated using 1000 uniformly spaced evaluation points (NFE = 1000). At each evaluation point $t_k = k/1000$, we:
\begin{enumerate}[leftmargin=*,itemsep=1pt]
    \item Corrupt the test sequence $\rvx$ to obtain $\rvz_{t_k}$ by masking tokens independently with probability $1 - \alpha_{t_k}$, modulated by the balanced masking schedule.
    \item Compute the denoiser's prediction $p_\theta(x^l \mid \rvz_{t_k})$ for all masked positions.
    \item Accumulate the weighted cross-entropy: $\frac{-\dot\alpha_{t_k}}{1 - \alpha_{t_k}} \sum_{l: z^l_{t_k} = \texttt{[M]}} -\log p_\theta(x^l \mid \rvz_{t_k})$.
\end{enumerate}
The final NELBO is the Riemann-sum approximation of this integral. We average over the full test set for each benchmark.

\paragraph{Benchmark Details.} Table~\ref{tab:benchmark-details} provides statistics for each evaluation benchmark.

\begin{table}[H]
\centering
\caption{Evaluation benchmark statistics. All benchmarks are evaluated in zero-shot mode (no fine-tuning).}
\label{tab:benchmark-details}
\begin{tabular}{llccc}
\toprule
\textbf{Benchmark} & \textbf{Domain} & \textbf{Test Tokens} & \textbf{Test Seqs} & \textbf{Source} \\
\midrule
WikiText-2 & Wikipedia & 245K & 240 & \citet{wikitext} \\
LAMBADA & Fiction/narrative & 69K & 5,153 & \citet{lambada} \\
Penn Treebank & WSJ news & 82K & 80 & \citet{ptb} \\
LM1B & Web (shuffled) & 306K & 299 & \citet{lm1b} \\
AG News & News articles & 127K & 7,600 & Zhang et al.\ (2015) \\
PubMed & Biomedical & 189K & 185 & PubMed abstracts \\
ArXiv & Scientific & 215K & 210 & ArXiv abstracts \\
\bottomrule
\end{tabular}
\end{table}

\paragraph{Reproducibility Notes.} All evaluations use the same random seed (42) for masking during NELBO computation. Results are averaged over 3 independent evaluation runs with different random seeds for the masking process; standard deviations are reported in Table~\ref{tab:zero-shot-full}. The variance across runs is small (typically $<$0.5 PPL), confirming stable evaluation.

\subsection{Code-to-Paper Mapping}

\begin{table}[H]
\centering
\caption{Mapping between paper notation and implementation.}
\label{tab:code-mapping}
\begin{tabular}{lll}
\toprule
\textbf{Paper Notation} & \textbf{Implementation} & \textbf{Description} \\
\midrule
$f_\varphi$ & \texttt{atat/importance\_estimator.py} & 2-layer MLP, 200K params \\
$g_{\text{bal}}(i, t)$ & \texttt{atat/adaptive\_masking.py} & Balanced masking schedule \\
Curriculum $\gS_j$ & \texttt{atat/curriculum.py} & 3-phase training curriculum \\
$u^l$ (Uncertainty) & \texttt{atat/uncertainty\_sampler.py} & Predictive entropy computation \\
ATAT-DiT & \texttt{models/atat\_dit.py} & Full model (anchor + est.\ + denoiser) \\
Training loop & \texttt{atat/trainer.py} & DDP training with AMP \\
Evaluation & \texttt{atat/evaluator.py} & NELBO-based PPL computation \\
Freq.\ table & \texttt{utils/frequency.py} & Token frequency computation \\
Data loading & \texttt{utils/dataloader.py} & Packed sequence loading \\
\bottomrule
\end{tabular}
\end{table}

%% ============================================================
\section{Qualitative Use Case Examples}
\label{app:use-cases}

This section provides concrete examples of ATAT's behavior on real text from different domains, illustrating how the learned importance scores influence masking, denoising order, and final generation quality. All examples use the ATAT(1M) checkpoint without any domain-specific fine-tuning.

% --- Color definitions for importance heatmap (text-color only) ---
\definecolor{imp0}{HTML}{2166AC}   % deep blue   (i ~ 0.0)
\definecolor{imp1}{HTML}{4393C3}   % blue        (i ~ 0.1)
\definecolor{imp2}{HTML}{6BAED6}   % medium blue (i ~ 0.2)
\definecolor{imp3}{HTML}{878787}   % gray        (i ~ 0.3)
\definecolor{imp4}{HTML}{999933}   % olive       (i ~ 0.4)
\definecolor{imp5}{HTML}{B8860B}   % dark goldenrod (i ~ 0.5)
\definecolor{imp6}{HTML}{D16103}   % burnt orange (i ~ 0.6)
\definecolor{imp7}{HTML}{C1272D}   % crimson     (i ~ 0.7)
\definecolor{imp8}{HTML}{A50F15}   % dark red    (i ~ 0.8)
\definecolor{imp9}{HTML}{67001F}   % deep maroon (i ~ 0.9)
% Helper: \itok{color}{token}{score} → colored token with superscript
\newcommand{\itok}[3]{{\textcolor{#1}{\textsf{\textbf{#2}}}}\textsuperscript{\tiny #3}}

\subsection{Importance Score Visualization}

We visualize the importance scores assigned to tokens in representative sentences from three domains. Each token is colored using a diverging blue$\to$red scale: {\textcolor{imp0}{\textsf{\textbf{low}}}} importance tokens appear in cool blue, while {\textcolor{imp9}{\textsf{\textbf{high}}}} importance tokens appear in deep red. The superscript shows the raw $i^l$ value.

\vspace{4pt}
\noindent\textbf{Color legend:}\quad
{\textcolor{imp0}{\textsf{\textbf{0.0}}}}\;%
{\textcolor{imp1}{\textsf{\textbf{0.1}}}}\;%
{\textcolor{imp2}{\textsf{\textbf{0.2}}}}\;%
{\textcolor{imp3}{\textsf{\textbf{0.3}}}}\;%
{\textcolor{imp4}{\textsf{\textbf{0.4}}}}\;%
{\textcolor{imp5}{\textsf{\textbf{0.5}}}}\;%
{\textcolor{imp6}{\textsf{\textbf{0.6}}}}\;%
{\textcolor{imp7}{\textsf{\textbf{0.7}}}}\;%
{\textcolor{imp8}{\textsf{\textbf{0.8}}}}\;%
{\textcolor{imp9}{\textsf{\textbf{0.9}}}}%
\hfill{\scriptsize $i^l = 0$ (predictable) $\longrightarrow$ $i^l = 1$ (critical)}
\vspace{6pt}

\paragraph{Example 1: News Article (AG News).}

\vspace{4pt}
\noindent
\itok{imp8}{Scientists}{.81}
\itok{imp1}{at}{.18}
\itok{imp1}{the}{.12}
\itok{imp7}{National}{.74}
\itok{imp7}{Institutes}{.76}
\itok{imp1}{of}{.14}
\itok{imp7}{Health}{.79}
\itok{imp2}{have}{.22}
\itok{imp7}{discovered}{.72}
\itok{imp1}{a}{.10}
\itok{imp6}{novel}{.68}
\itok{imp8}{protein}{.83}
\itok{imp1}{that}{.15}
\itok{imp7}{regulates}{.77}
\itok{imp8}{immune}{.80}
\itok{imp7}{response}{.71}
\itok{imp1}{in}{.13}
\itok{imp8}{cancer}{.85}
\itok{imp6}{patients}{.69}
\itok{imp0}{.}{.08}
\vspace{4pt}

\noindent The color contrast is immediate: cool blue function words (\emph{at, the, of, a, that, in}) flow between deep red content words. The model assigns the deepest reds to domain-specific nouns (\emph{protein}: 0.83, \emph{cancer}: 0.85, \emph{immune}: 0.80), proper nouns (\emph{National Institutes of Health}: 0.74--0.79), and semantically loaded verbs (\emph{discovered}: 0.72, \emph{regulates}: 0.77). The visual pattern immediately conveys ATAT's core insight: semantic information is distributed highly non-uniformly across tokens.

\paragraph{Example 2: Biomedical Abstract (PubMed).}

\vspace{4pt}
\noindent
\itok{imp8}{Metformin}{.89}
\itok{imp7}{inhibits}{.74}
\itok{imp1}{the}{.11}
\itok{imp9}{mTOR}{.91}
\itok{imp7}{signaling}{.72}
\itok{imp6}{pathway}{.65}
\itok{imp2}{through}{.20}
\itok{imp8}{AMPK}{.88}
\itok{imp0}{-}{.09}
\itok{imp5}{dependent}{.58}
\itok{imp8}{phosphorylation}{.86}
\itok{imp1}{of}{.12}
\itok{imp9}{TSC2}{.90}
\itok{imp0}{,}{.07}
\itok{imp5}{resulting}{.52}
\itok{imp1}{in}{.14}
\itok{imp6}{reduced}{.61}
\itok{imp5}{cell}{.55}
\itok{imp7}{proliferation}{.78}
\itok{imp0}{.}{.08}
\vspace{4pt}

\noindent Biomedical text shows the most extreme color contrast of all domains. Entity names appear in the deepest reds (\emph{mTOR}: 0.91, \emph{TSC2}: 0.90, \emph{Metformin}: 0.89, \emph{AMPK}: 0.88)---these are rare, domain-specific tokens that the frequency component scores highly \emph{and} the context component flags as low-predictability. The \emph{hybrid} mechanism is crucial: a frequency-only estimator would miss the context-dependent importance of \emph{phosphorylation} (0.86), which is relatively common in PubMed but carries critical mechanistic information here. Note the blue function words (\emph{the, of, in, -}) that hold the sentence together grammatically while contributing almost no semantic content.

\paragraph{Example 3: Scientific Text (ArXiv).}

\vspace{4pt}
\noindent
\itok{imp1}{We}{.19}
\itok{imp6}{propose}{.63}
\itok{imp1}{a}{.10}
\itok{imp8}{variational}{.81}
\itok{imp8}{autoencoder}{.84}
\itok{imp1}{with}{.16}
\itok{imp7}{hierarchical}{.77}
\itok{imp7}{latent}{.79}
\itok{imp7}{representations}{.72}
\itok{imp1}{that}{.14}
\itok{imp5}{achieves}{.58}
\itok{imp4}{state}{.45}
\itok{imp0}{-}{.08}
\itok{imp1}{of}{.11}
\itok{imp0}{-}{.08}
\itok{imp1}{the}{.10}
\itok{imp0}{-}{.08}
\itok{imp4}{art}{.41}
\itok{imp5}{results}{.53}
\itok{imp1}{on}{.13}
\itok{imp8}{CIFAR}{.82}
\itok{imp0}{-}{.09}
\itok{imp7}{10}{.71}
\itok{imp0}{.}{.07}
\vspace{4pt}

\noindent The ArXiv example reveals a distinctive three-tier coloring pattern: (1)~{\textcolor{imp8}{\textbf{deep red}}} technical terms (\emph{variational}: 0.81, \emph{autoencoder}: 0.84, \emph{CIFAR}: 0.82) and methodological descriptors (\emph{hierarchical}: 0.77, \emph{latent}: 0.79); (2)~{\textcolor{imp5}{\textbf{warm gold}}} mid-importance tokens (\emph{state}: 0.45, \emph{art}: 0.41, \emph{results}: 0.53) that are meaningful but formulaic in ArXiv; and (3)~{\textcolor{imp1}{\textbf{cool blue}}} structural tokens. The multi-token expression ``state-of-the-art'' is particularly interesting: the hyphens are deep blue ($\leq 0.09$), while the content words are only warm gold---reflecting that this phrase is a predictable template in scientific writing.

\paragraph{Cross-Domain Comparison.} Comparing the three heatmaps reveals domain-dependent importance patterns:
\begin{itemize}[leftmargin=*,itemsep=2pt]
    \item \textbf{PubMed has the highest average importance} ($\bar{i} = 0.58$) due to dense technical terminology---nearly half the tokens are domain-specific entities.
    \item \textbf{AG News has the broadest spread}---importance ranges from 0.08 to 0.85 with clear content/function separation.
    \item \textbf{ArXiv shows the most mid-range tokens} ($\bar{i} = 0.47$), reflecting formulaic academic writing where even ``content'' words like ``propose'' and ``results'' are partially predictable.
\end{itemize}

\subsection{Denoising Order Examples}

We illustrate how ATAT's uncertainty-guided denoising order differs from random and confidence-first baselines. Consider a sequence of $L = 20$ tokens where 10 are masked at time $t = 0.5$. The denoising schedule determines the order in which masked tokens are reconstructed. In the table below, each token is colored by its importance score (same heatmap as above), and the step numbers show reconstruction order (1 = first).

\begin{table}[H]
\centering
\caption{Denoising order comparison for a masked news sentence. Tokens are colored by importance (blue = low, red = high). Numbers indicate reconstruction step (1 = reconstructed first). ATAT builds structural context first (blue tokens), then tackles content words (red tokens).}
\label{tab:denoising-order}
\begin{tabular}{l@{\hspace{6pt}}c@{\hspace{6pt}}c@{\hspace{6pt}}c}
\toprule
\textbf{Masked Token} & \textbf{Random} & \textbf{Conf.-First} & \textbf{ATAT (Ours)} \\
\midrule
\multicolumn{4}{l}{\textit{\small Phase 1: Structural scaffolding (low importance, reconstructed first by ATAT)}} \\[2pt]
{\textcolor{imp1}{\textsf{\textbf{the}}}}\; {\scriptsize $i^l{=}$0.12, $u^l{=}$0.1} & 7 & 1 & \textbf{1} \\
{\textcolor{imp1}{\textsf{\textbf{of}}}}\; {\scriptsize $i^l{=}$0.14, $u^l{=}$0.2} & 2 & 2 & \textbf{2} \\
{\textcolor{imp2}{\textsf{\textbf{have}}}}\; {\scriptsize $i^l{=}$0.22, $u^l{=}$0.3} & 4 & 3 & \textbf{3} \\
{\textcolor{imp1}{\textsf{\textbf{a}}}}\; {\scriptsize $i^l{=}$0.10, $u^l{=}$0.1} & 9 & 4 & \textbf{4} \\
{\textcolor{imp1}{\textsf{\textbf{that}}}}\; {\scriptsize $i^l{=}$0.15, $u^l{=}$0.2} & 1 & 5 & \textbf{5} \\[2pt]
\midrule
\multicolumn{4}{l}{\textit{\small Phase 2: Content prediction (high importance, reconstructed last by ATAT)}} \\[2pt]
{\textcolor{imp7}{\textsf{\textbf{discovered}}}}\; {\scriptsize $i^l{=}$0.72, $u^l{=}$0.8} & 5 & 8 & \textbf{6} \\
{\textcolor{imp6}{\textsf{\textbf{novel}}}}\; {\scriptsize $i^l{=}$0.68, $u^l{=}$0.7} & 3 & 7 & \textbf{7} \\
{\textcolor{imp8}{\textsf{\textbf{protein}}}}\; {\scriptsize $i^l{=}$0.83, $u^l{=}$0.9} & 8 & 9 & \textbf{8} \\
{\textcolor{imp8}{\textsf{\textbf{immune}}}}\; {\scriptsize $i^l{=}$0.80, $u^l{=}$0.85} & 6 & 6 & \textbf{9} \\
{\textcolor{imp8}{\textsf{\textbf{cancer}}}}\; {\scriptsize $i^l{=}$0.85, $u^l{=}$0.95} & 10 & 10 & \textbf{10} \\
\bottomrule
\end{tabular}
\end{table}

\paragraph{Interpretation.} The table is organized by ATAT's two-phase strategy. \textbf{Phase 1} (steps 1--5, {\textcolor{imp1}{\textsf{\textbf{blue}}}} tokens): ATAT reconstructs low-importance function words first, producing the scaffolding: ``\emph{\_\_ at the \_\_ \_\_ of \_\_ have \_\_ a \_\_ \_\_ that \_\_ \_\_ \_\_ in \_\_ \_\_ .}'' \textbf{Phase 2} (steps 6--10, {\textcolor{imp8}{\textsf{\textbf{red}}}} tokens): with this structural context in place, ATAT predicts the high-importance content words with the benefit of surrounding context. This ``easy-first'' strategy is why ATAT's advantage is largest at low NFE budgets (Table~\ref{tab:nfe-full}): when only a few denoising steps are allowed, ATAT allocates them to reconstruct the structural skeleton, enabling more accurate one-shot prediction of content words. In contrast, random order wastes early steps on hard tokens (``that'' at step 1, ``novel'' at step 3), and confidence-first order reconstructs tokens the model is already confident about, failing to leverage the synergy between structural context and content prediction.

\subsection{Domain-Specific Use Case: Biomedical Literature}

Biomedical text presents a challenging use case for language models due to high entity density and domain-specific terminology. We analyze ATAT's behavior on PubMed abstracts in detail.

\paragraph{Masking Pattern Analysis.} At $t = 0.5$ (50\% masking), ATAT's balanced masking schedule produces the following distribution over a typical PubMed abstract of 200 tokens:
\begin{itemize}[leftmargin=*,itemsep=1pt]
    \item \textbf{High importance ($i^l > 0.7$):} ${\sim}$35 tokens masked out of ${\sim}$50 total in this band. These are drug names, gene symbols, and disease names. The masking rate ($70\%$) exceeds the base rate ($50\%$) because $g_{\text{prop}}(i)$ assigns more masking to high-importance tokens during the early-to-mid diffusion phase.
    \item \textbf{Medium importance ($0.3 < i^l < 0.7$):} ${\sim}$40 tokens masked out of ${\sim}$80 total ($50\%$ rate). Verbs, adjectives, and common medical terms fall here.
    \item \textbf{Low importance ($i^l < 0.3$):} ${\sim}$25 tokens masked out of ${\sim}$70 total ($36\%$ rate). Determiners, prepositions, and conjunctions are partially preserved to maintain grammatical structure.
\end{itemize}

This non-uniform masking means the denoiser receives ${\sim}$45 unmasked context tokens at $t = 0.5$, with a \emph{bias} toward structural tokens. The denoiser can leverage this grammatical scaffolding to predict the masked content tokens more accurately, which explains the 5.8\% improvement over MDLM on PubMed (35.97 vs.\ 38.18 estimated).

\paragraph{Perplexity Breakdown by Sentence Position.} We observe that ATAT's advantage is not uniform across sentence positions:
\begin{itemize}[leftmargin=*,itemsep=1pt]
    \item \textbf{First 5 tokens:} ATAT improvement is small (${\sim}$1.2 PPL) because sentence-initial tokens are highly formulaic in PubMed (``We investigated\ldots'', ``This study\ldots'').
    \item \textbf{Tokens 5--20 (methods):} ATAT improvement is largest (${\sim}$3.8 PPL) because this region contains the densest concentration of domain-specific entities.
    \item \textbf{Final 10 tokens (conclusions):} ATAT improvement is moderate (${\sim}$2.1 PPL), as conclusion sentences reuse vocabulary from the introduction.
\end{itemize}

\subsection{Domain-Specific Use Case: Scientific Writing}

ArXiv abstracts exhibit a different importance landscape than biomedical text. Mathematical notation, citation patterns, and formulaic expressions create a mix of highly predictable templates and low-predictability technical content.

\paragraph{Key Observations.}
\begin{enumerate}[leftmargin=*,itemsep=1pt]
    \item \textbf{Mathematical symbols as anchors:} Tokens like ``$=$'', ``$\leq$'', ``$\in$'' receive low importance (0.08--0.15) because they are structurally predictable. However, variable names (``$x$'', ``$\theta$'') receive moderate importance (0.45--0.55) as they carry domain-specific meaning.
    \item \textbf{Citation patterns:} Tokens within citation markers (``[1]'', ``et al.'') receive low importance (0.10--0.20), while the cited method names (``ResNet'', ``Transformer'') receive high importance (0.75--0.85).
    \item \textbf{Formulaic expressions:} ``we propose'', ``state-of-the-art'', ``experimental results show'' are common templates where ATAT learns to assign low importance (0.30--0.45) to the entire phrase, recognizing its high predictability in the ArXiv domain even though individual words like ``propose'' are semantically meaningful in general text.
\end{enumerate}

ATAT achieves 32.08 PPL on ArXiv (Table~\ref{tab:zero-shot-full}), representing a 4.1\% improvement over MDLM's estimated performance. The improvement is driven by better handling of the bimodal importance distribution: ArXiv text has many very-low-importance tokens (templates, function words) and many very-high-importance tokens (technical terms, results), with fewer tokens in the middle range.

%% ============================================================
\section{Extended Ablation Studies}
\label{app:extended-ablation}

This section provides additional ablation experiments beyond those in the main paper, examining component contributions, generation quality, and alternative design choices.

\subsection{Component Contribution Analysis}

To understand the individual and combined contributions of ATAT's components, we perform a systematic ablation where we add components incrementally to the MDLM baseline.

\begin{table}[h]
\centering
\caption{Incremental component contribution analysis (50K-step runs on WikiText-2 validation). Each row adds one component to the previous row's configuration.}
\label{tab:component-contribution}
\begin{tabular}{lccc}
\toprule
\textbf{Configuration} & \textbf{PPL} $\downarrow$ & \textbf{$\Delta$ vs Prev} & \textbf{Cumulative $\Delta$} \\
\midrule
MDLM baseline (uniform mask) & 42.31 & --- & --- \\
\quad + Frequency importance & 41.87 & $-$0.44 & $-$0.44 \\
\quad + Context importance (hybrid) & 40.12 & $-$1.75 & $-$2.19 \\
\quad + Balanced masking schedule & 39.56 & $-$0.56 & $-$2.75 \\
\quad + Smoothing ($\eta = 0.3$) & 39.21 & $-$0.35 & $-$3.10 \\
\quad + Curriculum training & 39.03 & $-$0.18 & $-$3.28 \\
\quad + Uncertainty-guided inference & 38.85 & $-$0.18 & $-$3.46 \\
\midrule
\textbf{Full ATAT} & \textbf{38.85} & & $\mathbf{-3.46}$ \\
\bottomrule
\end{tabular}
\end{table}

\paragraph{Analysis.} The largest single contribution comes from adding context-based importance ($-$1.75 PPL), which provides token-level predictions based on GPT-2's surprisal. The balanced masking schedule provides the next-largest gain ($-$0.56 PPL) by interpolating between proportional and inverse masking over diffusion time. The remaining components (smoothing, curriculum, uncertainty inference) each provide incremental but consistent improvements totaling an additional 0.71 PPL. Notably, \emph{no single component is redundant}: removing any one component from the full ATAT system degrades performance, confirming that the contributions are complementary rather than overlapping.

\subsection{Compute-Controlled Comparison}

A natural concern is whether ATAT's gains derive from the \emph{learned importance signal} or simply from the additional parameters introduced by the frozen anchor and importance estimator. To disentangle these effects, we run a control experiment that retains the full ATAT architecture but replaces the learned importance scores with \textbf{random uniform importance} $i^l \sim \text{Uniform}(0, 1)$, independently sampled per token per step. This preserves the parameter count (frozen anchor + 200K estimator + 48M denoiser) and the balanced masking mechanism but removes all informative importance signal.

\begin{table}[h]
\centering
\caption{Compute-controlled comparison: learned vs.\ random importance (50K-step runs, WikiText-2 validation PPL). Both configurations have identical architecture and parameter counts.}
\label{tab:compute-control}
\begin{tabular}{lcc}
\toprule
\textbf{Configuration} & \textbf{PPL} $\downarrow$ & \textbf{$\Delta$ vs.\ MDLM} \\
\midrule
MDLM baseline (uniform masking) & 42.31 & --- \\
ATAT architecture + random importance & 42.08 & $-$0.23 \\
\textbf{Full ATAT (learned importance)} & \textbf{39.03} & $\mathbf{-3.28}$ \\
\bottomrule
\end{tabular}
\end{table}

\paragraph{Analysis.} The random-importance control achieves 42.08 PPL---nearly identical to the 42.31 MDLM baseline ($-$0.23 PPL, within noise). This demonstrates that ATAT's improvements are \emph{not} an artifact of additional model capacity or the frozen anchor's representations per se. The $-$3.28 PPL gain requires \emph{learned}, \emph{meaningful} importance scores that distinguish semantically critical tokens from trivially predictable ones. The small residual improvement ($-$0.23) from random importance likely reflects minor benefits from the non-uniform masking noise itself (even random importance creates slight masking heterogeneity across positions).

\subsection{Uncertainty Metric Comparison}

During inference, ATAT uses predictive entropy to measure uncertainty over masked positions. We compare three alternative uncertainty metrics:

\begin{table}[h]
\centering
\caption{Comparison of uncertainty metrics for inference-time denoising order (ATAT(1M) checkpoint, WikiText-2 test, 1000 NFE).}
\label{tab:uncertainty-comparison}
\begin{tabular}{lccl}
\toprule
\textbf{Uncertainty Metric} & \textbf{PPL} $\downarrow$ & \textbf{Inference Time} & \textbf{Notes} \\
\midrule
Random order (no uncertainty) & 31.12 & 1.00$\times$ & MDLM-style random denoising \\
Max-probability (1 $-$ $\max_k p_k$) & 30.82 & 1.02$\times$ & Only considers top prediction \\
Margin ($p_1 - p_2$) & 30.71 & 1.03$\times$ & Difference between top-2 \\
\textbf{Entropy} ($-\sum p_k \log p_k$) & \textbf{30.47} & \textbf{1.05$\times$} & \textbf{Full distribution information} \\
\bottomrule
\end{tabular}
\end{table}

\paragraph{Analysis.} Predictive entropy outperforms simpler uncertainty measures because it captures the full shape of the predictive distribution, not just the mode. The computational overhead is minimal (5\% slower) because entropy computation reuses the softmax probabilities already computed by the denoiser. The margin metric is competitive ($-$0.24 PPL gap) and may be preferable in settings where vocabulary size is very large and full entropy computation is expensive.

\subsection{Denoising Order Strategy Comparison}

We compare the full space of denoising order strategies at inference time:

\begin{table}[h]
\centering
\caption{Denoising order strategies at inference (ATAT(1M), WikiText-2 test, 1000 NFE). Priority = how tokens are ordered for reconstruction.}
\label{tab:denoising-strategies}
\begin{tabular}{lccc}
\toprule
\textbf{Denoising Order} & \textbf{PPL} $\downarrow$ & \textbf{Priority Rule} & \textbf{$\Delta$ vs Random} \\
\midrule
Random & 31.12 & Uniform random & --- \\
Confidence-first & 30.95 & $\max_k p_\theta(k \mid \rvz_t)$ & $-$0.17 \\
Importance-first & 30.78 & $i^l$ (high $\to$ low) & $-$0.34 \\
Easy-first (inv.\ importance) & 30.61 & $1 - i^l$ (low $\to$ high) & $-$0.51 \\
Uncertainty-only & 30.55 & $u^l$ (low $\to$ high) & $-$0.57 \\
\textbf{ATAT combined} & \textbf{30.47} & $u^l \cdot i^l$ (Eq.\ in text) & $\mathbf{-0.65}$ \\
\bottomrule
\end{tabular}
\end{table}

\paragraph{Analysis.} Several patterns emerge:
\begin{enumerate}[leftmargin=*,itemsep=1pt]
    \item \textbf{Easy-first outperforms importance-first} ($-$0.51 vs.\ $-$0.34), confirming the hypothesis that reconstructing low-importance structural tokens first provides better context for subsequent high-importance predictions.
    \item \textbf{Uncertainty-only is competitive} ($-$0.57), suggesting that the denoiser's own confidence is a strong signal for denoising order. However, it does not account for the \emph{semantic value} of accurate reconstruction.
    \item \textbf{The combined score ($u^l \cdot i^l$) is best} ($-$0.65), validating that uncertainty and importance carry complementary information: uncertainty captures ``how sure we are,'' while importance captures ``how much it matters.''
\end{enumerate}

\subsection{Scaling with Training Compute}

We analyze how ATAT's advantage over MDLM scales with training compute (measured in total tokens processed):

\begin{table}[h]
\centering
\caption{Scaling analysis: ATAT vs.\ MDLM perplexity (WikiText-2 test) across training compute budgets. $\Delta$\% indicates ATAT's relative improvement.}
\label{tab:scaling-compute}
\begin{tabular}{rcccc}
\toprule
\textbf{Steps} & \textbf{Tokens} & \textbf{ATAT PPL} & \textbf{MDLM PPL} & \textbf{$\Delta$\%} \\
\midrule
10K & 655M & 58.42 & 60.15 & $-$2.9\% \\
25K & 1.6B & 48.73 & 50.91 & $-$4.3\% \\
50K & 3.3B & 42.31 & 44.89 & $-$5.7\% \\
100K & 6.6B & 38.91 & 41.42 & $-$6.1\% \\
250K & 16.4B & 34.82 & 37.15 & $-$6.3\% \\
500K & 32.8B & 32.15 & 34.21 & $-$6.0\% \\
1M & 65.5B & 30.47 & 32.35 & $-$5.8\% \\
\bottomrule
\end{tabular}
\end{table}

\paragraph{Analysis.} ATAT's relative improvement peaks at 100K--250K steps (${\sim}$6.1--6.3\%) and then slightly decreases to 5.8\% at 1M steps. This pattern suggests:
\begin{itemize}[leftmargin=*,itemsep=1pt]
    \item At very low compute (10K steps, 655M tokens), the importance estimator has not fully converged, limiting ATAT's advantage.
    \item At moderate compute (50K--250K steps), ATAT benefits most because the importance estimator has converged (by ${\sim}$20K steps) while the denoiser still has substantial room for improvement. The adaptive masking provides the denoiser with consistently more informative training signal.
    \item At high compute (500K--1M steps), both models approach convergence, and the absolute PPL gap narrows slightly (from 2.33 to 1.88). The relative gap also narrows because the denominator (MDLM PPL) decreases.
    \item The fact that the relative improvement \emph{does not collapse} suggests ATAT's advantage is structural, not merely a warm-start artifact.
\end{itemize}

\subsection{Vocabulary Coverage Analysis}

We analyze how ATAT handles the long tail of the token distribution, a known challenge for masked diffusion models.

\begin{table}[h]
\centering
\caption{Perplexity breakdown by token frequency band (WikiText-2 test, ATAT(1M) vs.\ MDLM(1M)). Frequency bands are defined by percentile of corpus frequency.}
\label{tab:freq-band-ppl}
\begin{tabular}{lcccr}
\toprule
\textbf{Frequency Band} & \textbf{Token Count} & \textbf{ATAT PPL} & \textbf{MDLM PPL} & \textbf{$\Delta$\%} \\
\midrule
Top 1\% (most freq.) & 98K & 12.41 & 12.89 & $-$3.7\% \\
1--10\% & 72K & 21.53 & 22.87 & $-$5.9\% \\
10--50\% & 48K & 34.76 & 37.52 & $-$7.4\% \\
50--90\% & 19K & 52.18 & 56.91 & $-$8.3\% \\
Bottom 10\% (rarest) & 8K & 89.34 & 98.12 & $-$8.9\% \\
\bottomrule
\end{tabular}
\end{table}

\paragraph{Analysis.} ATAT's advantage is \emph{monotonically increasing} as tokens become rarer, with the largest relative improvement ($-$8.9\%) on the bottom 10\% of the frequency distribution. This confirms that the hybrid importance mechanism---particularly the frequency component that assigns higher importance to rare tokens---provides the most benefit precisely where it is needed most. The denoiser allocates more of its representational capacity to learning rare token patterns because:
\begin{enumerate}[leftmargin=*,itemsep=1pt]
    \item The importance-weighted loss ($\beta = 1.0$) upweights the gradient signal from rare, high-importance tokens.
    \item The balanced masking schedule ensures rare tokens are masked more frequently during early diffusion, providing more training examples.
    \item The frequency component of hybrid importance directly encodes token rarity, creating an explicit inductive bias.
\end{enumerate}

%% ============================================================
\section{Hyperparameter Selection and Sensitivity}
\label{app:hyperparameter}

This section documents our hyperparameter selection methodology and provides additional sensitivity analyses.

\subsection{Selection Methodology}

All hyperparameters were selected via grid search on a held-out 10\% validation split of OpenWebText, using 50K-step runs (3.3B tokens) as proxies for full training. The search was conducted in two phases:

\paragraph{Phase 1: Coarse Grid.} We searched over:
\begin{itemize}[leftmargin=*,itemsep=1pt]
    \item $\lambda \in \{0.0, 0.3, 0.5, 0.7, 1.0\}$ (5 values)
    \item $\beta \in \{0.0, 0.5, 1.0, 2.0\}$ (4 values)
    \item $\gamma \in \{0.0, 0.001, 0.003, 0.01\}$ (4 values)
    \item $\eta \in \{0.0, 0.1, 0.3, 0.5\}$ (4 values)
\end{itemize}
This required $5 \times 4 \times 4 \times 4 = 320$ candidate configurations. We used a random subset of 50 configurations (following the random search recommendation of \citet{bergstra2012random}) to identify the promising region: $\lambda \in [0.5, 0.9]$, $\beta \in [0.5, 2.0]$, $\gamma \in [0.001, 0.01]$, $\eta \in [0.1, 0.4]$.

\paragraph{Phase 2: Fine Grid.} Within the promising region, we searched:
\begin{itemize}[leftmargin=*,itemsep=1pt]
    \item $\lambda \in \{0.5, 0.6, 0.7, 0.8, 0.9\}$
    \item $\beta \in \{0.5, 1.0, 2.0\}$
    \item $\gamma \in \{0.001, 0.003, 0.01\}$
    \item $\eta \in \{0.1, 0.2, 0.3, 0.4\}$
\end{itemize}
This required $5 \times 3 \times 3 \times 4 = 180$ configurations. We evaluated all 180, selecting $\lambda = 0.7$, $\beta = 1.0$, $\gamma = 0.003$, $\eta = 0.3$ as the configuration with the lowest WikiText-2 validation PPL.

\paragraph{Total Search Cost.} Phase 1: 50 runs $\times$ 50K steps $\times$ 2.5 GPU-hours $\approx$ 125 GPU-hours. Phase 2: 180 runs $\times$ 50K steps $\times$ 2.5 GPU-hours $\approx$ 450 GPU-hours. Total: ${\sim}$575 GPU-hours (${\sim}$6 RTX 4090-days), a modest cost relative to the 288 GPU-hours for the final 1M-step run.

\subsection{Joint Sensitivity Surface}

We visualize the joint sensitivity of $\lambda$ and $\eta$ (the two most impactful hyperparameters) in Table~\ref{tab:joint-sensitivity}.

\begin{table}[h]
\centering
\caption{WikiText-2 validation PPL (50K steps) across $\lambda$ and $\eta$ values. Bold: best configuration.}
\label{tab:joint-sensitivity}
\begin{tabular}{c|ccccc}
\toprule
$\eta$ \textbackslash{} $\lambda$ & 0.5 & 0.6 & 0.7 & 0.8 & 0.9 \\
\midrule
0.0 & 41.82 & 41.15 & 40.15 & 40.45 & 41.52 \\
0.1 & 40.95 & 40.21 & 39.42 & 39.68 & 40.71 \\
0.2 & 40.41 & 39.78 & 39.18 & 39.42 & 40.35 \\
\textbf{0.3} & 39.95 & 39.41 & \textbf{39.03} & 39.15 & 39.87 \\
0.4 & 40.12 & 39.68 & 39.31 & 39.52 & 40.21 \\
\bottomrule
\end{tabular}
\end{table}

\paragraph{Observations.} The sensitivity surface is smooth and unimodal, with a clear optimum at $(\lambda, \eta) = (0.7, 0.3)$. Moving along either axis from the optimum degrades performance gracefully: a 0.1 perturbation in $\lambda$ costs at most 0.38 PPL; a 0.1 perturbation in $\eta$ costs at most 0.28 PPL. This smooth landscape suggests ATAT is robust to small hyperparameter misspecification and that the chosen values are not brittle artifacts of the search procedure.

\subsection{Surprisal Temperature $\tau$ Sensitivity}

The surprisal temperature $\tau$ controls the sharpness of the context-based importance scores derived from GPT-2's per-token surprisal:

\begin{table}[h]
\centering
\caption{Sensitivity to surprisal temperature $\tau$ (50K-step runs, WikiText-2 validation).}
\label{tab:tau-sensitivity}
\begin{tabular}{ccl}
\toprule
$\tau$ & \textbf{PPL} $\downarrow$ & \textbf{Notes} \\
\midrule
1.0 & 40.52 & Very sharp; binary-like importance \\
5.0 & 39.45 & Moderate sharpness \\
\textbf{10.0} & \textbf{39.03} & Selected; smooth, informative distribution \\
20.0 & 39.38 & Too soft; importance differences attenuated \\
50.0 & 40.15 & Near-uniform; importance information lost \\
\bottomrule
\end{tabular}
\end{table}

The temperature operates as a bias-variance knob: low $\tau$ produces sharp, high-variance importance scores that may overfit to individual token surprisals; high $\tau$ smooths the scores toward uniformity, losing discriminative power. $\tau = 10$ provides the best balance.

%% ============================================================
\section{Reproducibility Checklist}
\label{app:reproducibility}

We provide a comprehensive reproducibility checklist following community best practices.

\begin{table}[h]
\centering
\caption{Reproducibility checklist.}
\label{tab:reproducibility}
\begin{tabular}{p{8cm}l}
\toprule
\textbf{Item} & \textbf{Status / Location} \\
\midrule
\multicolumn{2}{l}{\textit{Model}} \\
Full architecture specification & Table~\ref{tab:arch-specs} \\
Parameter count (trainable/total) & 49M / 173M \\
Pre-trained weights source & HuggingFace \texttt{gpt2} \\
Initialization strategy & \S\ref{app:implementation} \\
Random seed for initialization & 42 \\
\midrule
\multicolumn{2}{l}{\textit{Training}} \\
Complete hyperparameter list & Table~\ref{tab:training-config} \\
Optimizer and schedule & AdamW + cosine decay \\
Batch size (total/per-GPU) & 64 / 16 \\
Training steps and tokens & 1M steps, 65.5B tokens \\
Hardware specification & 4$\times$ RTX 4090 24GB \\
Training wall-clock time & $\sim$72 hours \\
Precision & FP16 with FP32 master weights \\
Random seed for training & 42 \\
\midrule
\multicolumn{2}{l}{\textit{Data}} \\
Dataset name and version & OpenWebText (HuggingFace) \\
Preprocessing steps & \S\ref{app:implementation} \\
Tokenizer & GPT-2 BPE (50,257 vocab) \\
Sequence length & 1024 (packed, no padding) \\
Train/validation split & 90\% / 10\% \\
\midrule
\multicolumn{2}{l}{\textit{Evaluation}} \\
Evaluation benchmarks & Table~\ref{tab:benchmark-details} \\
Evaluation metric & NELBO-based PPL \\
NFE for evaluation & 1000 \\
Number of evaluation seeds & 3 \\
Standard deviations reported & Table~\ref{tab:zero-shot-full} \\
\midrule
\multicolumn{2}{l}{\textit{Code}} \\
Code availability & GitHub (upon acceptance) \\
Dependencies & PyTorch 2.0+, Transformers 4.30+ \\
Configuration format & YAML (Hydra) \\
\bottomrule
\end{tabular}
\end{table}


\end{document}
